% Chapter 7 — weave the graph: one build, every representation (kym_kg).

\gmzoomin{graph}
\note{Chapter 7: everything we learned becomes one graph.}

\begin{frame}{Doge, step 6: everything we learned, as dots and arrows}
  \centering
  \resizebox{\linewidth}{!}{% generated by scripts/figures.py — do not edit
\begin{tikzpicture}[x=1cm, y=1cm,
 dot/.style={circle, inner sep=0pt, minimum size=#1},
 lab/.style={font=\fontsize{7.4}{8.4}\selectfont, text=gmInk, inner sep=.8pt},
 edge/.style={draw=gmBase, line width=.35pt}]
\node[dot=8mm, fill=kFrame] (doge) at (0.00,0.00) {};
\node[font=\scriptsize\bfseries, text=white] at (doge) {Doge};
\node[dot=3.2mm, fill=kFrame] (anc0) at (0.00,1.00) {};
\node[lab, anchor=west] at (anc0.east) {Interior Monologue Captioning};
\node[dot=3.2mm, fill=kFrame] (anc1) at (0.00,1.62) {};
\node[lab, anchor=west] at (anc1.east) {Image Macros};
\node[dot=3.2mm, fill=kFrame] (anc2) at (0.00,2.24) {};
\node[lab, anchor=west] at (anc2.east) {Memes};
\node[dot=2.8mm, fill=kFrame] (ch0) at (-1.70,2.60) {};
\node[lab, anchor=east] at (ch0.west) {Doge 2 / Caesar};
\node[dot=2.8mm, fill=kFrame] (ch1) at (-2.00,2.17) {};
\node[lab, anchor=east] at (ch1.west) {Ironic Doge Memes};
\node[dot=2.8mm, fill=kFrame] (ch2) at (-1.70,1.73) {};
\node[lab, anchor=east] at (ch2.west) {Shiba Inus / Shibes};
\node[dot=2.8mm, fill=kFrame] (ch3) at (-2.00,1.30) {};
\node[lab, anchor=east] at (ch3.west) {The Death Of Kabosu / Doge};
\node[dot=2.3mm, fill=kFrame] (sb0) at (-5.10,1.00) {};
\node[lab, anchor=east] at (sb0.west) {Oh You So Crazy};
\node[dot=2.3mm, fill=kFrame] (sb1) at (-5.40,0.68) {};
\node[lab, anchor=east] at (sb1.west) {Stop Das Gay};
\node[dot=2.3mm, fill=kFrame] (sb2) at (-5.10,0.36) {};
\node[lab, anchor=east] at (sb2.west) {Animals Talking In All Caps};
\node[dot=2.3mm, fill=kFrame] (sb3) at (-5.40,0.04) {};
\node[lab, anchor=east] at (sb3.west) {Laughing Lizard / hhhehehe};
\node[dot=2.3mm, fill=kFrame] (sb4) at (-5.10,-0.29) {};
\node[lab, anchor=east] at (sb4.west) {Captioned Stock Photos};
\node[dot=2.3mm, fill=kFrame] (sb5) at (-5.40,-0.61) {};
\node[lab, anchor=east] at (sb5.west) {Master Trole Kid};
\node[dot=2.3mm, fill=kFrame] (sb6) at (-5.10,-0.93) {};
\node[lab, anchor=east] at (sb6.west) {Snek};
\node[dot=2.3mm, fill=kFrame] (sb7) at (-5.40,-1.25) {};
\node[lab, anchor=east] at (sb7.west) {Stop It Son, You Are Doing Me A Frighten};
\node[dot=2.2mm, fill=kWikidata] (wdQ384060) at (3.80,2.55) {};
\node[lab, anchor=west] at (wdQ384060.east) {Tumblr};
\node[dot=2.2mm, fill=kWikidata] (wdQ39315) at (3.80,1.99) {};
\node[lab, anchor=west] at (wdQ39315.east) {Shiba Inu};
\node[dot=2.2mm, fill=kWikidata] (wdQ23486479) at (3.80,1.43) {};
\node[lab, anchor=west] at (wdQ23486479.east) {Kabosu};
\node[dot=2.2mm, fill=kWikidata] (wdQ15613810) at (3.80,0.87) {};
\node[lab, anchor=west] at (wdQ15613810.east) {Doge};
\node[dot=2.2mm, fill=kWikidata] (wdQ261197) at (3.80,0.31) {};
\node[lab, anchor=west] at (wdQ261197.east) {monologue};
\node[dot=2.2mm, fill=kWikidata] (wdQ868492) at (3.80,-0.25) {};
\node[lab, anchor=west] at (wdQ868492.east) {Comic Sans MS};
\node[dot=2.2mm, fill=kWikidata] (wdQ144) at (5.85,2.55) {};
\node[lab, anchor=west] at (wdQ144.east) {dog};
\node[dot=2.2mm, fill=kWikidata] (wdQ131723) at (5.85,1.99) {};
\node[lab, anchor=west] at (wdQ131723.east) {Bitcoin};
\node[dot=2.2mm, fill=kWikidata] (wdQ1136) at (5.85,1.43) {};
\node[lab, anchor=west] at (wdQ1136.east) {Reddit};
\node[dot=2.2mm, fill=kWikidata] (wdQ15377916) at (5.85,0.87) {};
\node[lab, anchor=west] at (wdQ15377916.east) {Dogecoin};
\node[dot=2.2mm, fill=kWikidata] (wdQ5287) at (5.85,0.31) {};
\node[lab, anchor=west] at (wdQ5287.east) {Japanese};
\node[dot=2.7mm, fill=kTemplate] (t0) at (2.45,-0.75) {};
\node[lab, anchor=west] at (t0.east) {Doge};
\node[dot=2.7mm, fill=kTemplate] (t1) at (2.45,-1.06) {};
\node[lab, anchor=west] at (t1.east) {Average Fan vs. Average Enjoyer};
\node[dot=2.7mm, fill=kTemplate] (t2) at (2.45,-1.38) {};
\node[lab, anchor=west] at (t2.east) {Doge 2};
\node[dot=2.7mm, fill=kTemplate] (t3) at (2.45,-1.69) {};
\node[lab, anchor=west] at (t3.east) {doge then and now};
\node[dot=2.7mm, fill=kTemplate] (t4) at (2.45,-2.01) {};
\node[lab, anchor=west] at (t4.east) {Dogecoin};
\node[dot=2.7mm, fill=kTemplate] (t5) at (2.45,-2.32) {};
\node[lab, anchor=west] at (t5.east) {Swole Doge vs Cheems};
\node[dot=2.7mm, fill=kTemplate] (t6) at (2.45,-2.64) {};
\node[lab, anchor=west] at (t6.east) {Buff Doge vs. Cheems};
\node[dot=2.7mm, fill=kTemplate] (t7) at (2.45,-2.95) {};
\node[lab, anchor=west] at (t7.east) {Multi Doge};
\node[dot=2.2mm, fill=kWikidata] (wdQ33767) at (6.55,-1.10) {};
\node[lab, anchor=west] at (wdQ33767.east) {hand};
\node[dot=2.2mm, fill=kWikidata] (wdQ131514) at (6.55,-1.63) {};
\node[lab, anchor=west] at (wdQ131514.east) {couch};
\node[dot=2.2mm, fill=kWikidata] (wdQ317521) at (6.55,-2.17) {};
\node[lab, anchor=west] at (wdQ317521.east) {Elon Musk};
\node[dot=2.2mm, fill=kWikidata] (wdQ110983191) at (6.55,-2.70) {};
\node[lab, anchor=west] at (wdQ110983191.east) {Balltze};
\node[dot=2.5mm, fill=kConcept] (ty0) at (-1.60,-1.55) {};
\node[lab, anchor=east] at (ty0.west) {slang};
\node[dot=2.5mm, fill=kConcept] (ty1) at (-1.60,-1.81) {};
\node[lab, anchor=east] at (ty1.west) {character};
\node[dot=2.5mm, fill=kConcept] (ty2) at (-1.60,-2.08) {};
\node[lab, anchor=east] at (ty2.west) {animal};
\node[dot=2.5mm, fill=kConcept] (ty3) at (-1.60,-2.34) {};
\node[lab, anchor=east] at (ty3.west) {image-macro};
\node[dot=2.5mm, fill=kConcept] (ty4) at (-1.60,-2.60) {};
\node[lab, anchor=east] at (ty4.west) {exploitable};
\node[dot=1.6mm, fill=kConcept] (tg0) at (-3.60,-1.55) {};
\node[dot=1.6mm, fill=kConcept] (tg1) at (-3.60,-1.91) {};
\node[dot=1.6mm, fill=kConcept] (tg2) at (-3.60,-2.27) {};
\node[dot=1.6mm, fill=kConcept] (tg3) at (-3.60,-2.63) {};
\node[dot=1.6mm, fill=kConcept] (tg4) at (-4.02,-1.55) {};
\node[dot=1.6mm, fill=kConcept] (tg5) at (-4.02,-1.91) {};
\node[dot=1.6mm, fill=kConcept] (tg6) at (-4.02,-2.27) {};
\node[dot=1.6mm, fill=kConcept] (tg7) at (-4.02,-2.63) {};
\node[dot=1.6mm, fill=kConcept] (tg8) at (-4.44,-1.55) {};
\node[dot=1.6mm, fill=kConcept] (tg9) at (-4.44,-1.91) {};
\node[dot=1.6mm, fill=kConcept] (tg10) at (-4.44,-2.27) {};
\node[dot=1.6mm, fill=kConcept] (tg11) at (-4.44,-2.63) {};
\node[dot=1.6mm, fill=kConcept] (tg12) at (-4.86,-1.55) {};
\node[dot=1.6mm, fill=kConcept] (tg13) at (-4.86,-1.91) {};
\node[dot=1.6mm, fill=kConcept] (tg14) at (-4.86,-2.27) {};
\node[dot=1.6mm, fill=kConcept] (tg15) at (-4.86,-2.63) {};
\node[dot=1.6mm, fill=kConcept] (tg16) at (-5.28,-1.55) {};
\node[dot=1.6mm, fill=kConcept] (tg17) at (-5.28,-1.91) {};
\node[dot=1.6mm, fill=kConcept] (tg18) at (-5.28,-2.27) {};
\node[dot=1.6mm, fill=kConcept] (tg19) at (-5.28,-2.63) {};
\node[dot=1.6mm, fill=kConcept] (tg20) at (-5.70,-1.55) {};
\node[lab, text=gmInk2, anchor=north] at (-4.65,-2.95) {21 tags};
\node[dot=2.0mm, fill=kEvent] (ev0) at (-1.05,-2.30) {};
\node[dot=2.0mm, fill=kEvent] (ev1) at (-0.91,-2.34) {};
\node[dot=2.0mm, fill=kEvent] (ev2) at (-0.77,-2.37) {};
\node[dot=2.0mm, fill=kEvent] (ev3) at (-0.63,-2.41) {};
\node[dot=2.0mm, fill=kEvent] (ev4) at (-0.49,-2.43) {};
\node[dot=2.0mm, fill=kEvent] (ev5) at (-0.35,-2.46) {};
\node[dot=2.0mm, fill=kEvent] (ev6) at (-0.21,-2.47) {};
\node[dot=2.0mm, fill=kEvent] (ev7) at (-0.07,-2.48) {};
\node[dot=2.0mm, fill=kEvent] (ev8) at (0.07,-2.48) {};
\node[dot=2.0mm, fill=kEvent] (ev9) at (0.21,-2.47) {};
\node[dot=2.0mm, fill=kEvent] (ev10) at (0.35,-2.46) {};
\node[dot=2.0mm, fill=kEvent] (ev11) at (0.49,-2.43) {};
\node[dot=2.0mm, fill=kEvent] (ev12) at (0.63,-2.41) {};
\node[dot=2.0mm, fill=kEvent] (ev13) at (0.77,-2.37) {};
\node[dot=2.0mm, fill=kEvent] (ev14) at (0.91,-2.34) {};
\node[dot=2.0mm, fill=kEvent] (ev15) at (1.05,-2.30) {};
\node[lab, text=gmInk2, anchor=north] at (0,-2.62) {16 events, one story};
\begin{pgfonlayer}{background}
\draw[edge, draw=kFrame!60!gmSurface, line width=.7pt, -{Stealth[length=1.5mm]}] (doge) -- (anc0);
\draw[edge, draw=kFrame!60!gmSurface, line width=.7pt, -{Stealth[length=1.5mm]}] (anc0) -- (anc1);
\draw[edge, draw=kFrame!60!gmSurface, line width=.7pt, -{Stealth[length=1.5mm]}] (anc1) -- (anc2);
\draw[edge, draw=kFrame!60!gmSurface, -{Stealth[length=1.4mm]}] (ch0) -- (doge);
\draw[edge, draw=kFrame!60!gmSurface, -{Stealth[length=1.4mm]}] (ch1) -- (doge);
\draw[edge, draw=kFrame!60!gmSurface, -{Stealth[length=1.4mm]}] (ch2) -- (doge);
\draw[edge, draw=kFrame!60!gmSurface, -{Stealth[length=1.4mm]}] (ch3) -- (doge);
\draw[edge, dashed, draw=kFrame!45!gmSurface] (doge) -- (sb0);
\draw[edge, dashed, draw=kFrame!45!gmSurface] (doge) -- (sb1);
\draw[edge, dashed, draw=kFrame!45!gmSurface] (doge) -- (sb2);
\draw[edge, dashed, draw=kFrame!45!gmSurface] (doge) -- (sb3);
\draw[edge, dashed, draw=kFrame!45!gmSurface] (doge) -- (sb4);
\draw[edge, dashed, draw=kFrame!45!gmSurface] (doge) -- (sb5);
\draw[edge, dashed, draw=kFrame!45!gmSurface] (doge) -- (sb6);
\draw[edge, dashed, draw=kFrame!45!gmSurface] (doge) -- (sb7);
\draw[edge] (doge) -- (wdQ384060);
\draw[edge] (doge) -- (wdQ39315);
\draw[edge] (doge) -- (wdQ23486479);
\draw[edge] (doge) -- (wdQ15613810);
\draw[edge] (doge) -- (wdQ261197);
\draw[edge] (doge) -- (wdQ868492);
\draw[edge] (doge) -- (wdQ144);
\draw[edge] (doge) -- (wdQ131723);
\draw[edge] (doge) -- (wdQ1136);
\draw[edge] (doge) -- (wdQ15377916);
\draw[edge] (doge) -- (wdQ5287);
\draw[edge] (doge) -- (t0);
\draw[edge] (doge) -- (t1);
\draw[edge] (doge) -- (t2);
\draw[edge] (doge) -- (t3);
\draw[edge] (doge) -- (t4);
\draw[edge] (doge) -- (t5);
\draw[edge] (doge) -- (t6);
\draw[edge] (doge) -- (t7);
\draw[edge, draw=kTemplate!55!gmSurface] (t2) -- (wdQ33767);
\draw[edge, draw=kTemplate!55!gmSurface] (t2) -- (wdQ131514);
\draw[edge, draw=kTemplate!55!gmSurface] (t4) -- (wdQ317521);
\draw[edge, draw=kTemplate!55!gmSurface] (t5) -- (wdQ110983191);
\draw[edge, draw=kTemplate!55!gmSurface] (t6) -- (wdQ110983191);
\draw[edge, draw=kConcept!40!gmSurface] (doge) -- (ty0);
\draw[edge, draw=kConcept!40!gmSurface] (doge) -- (ty1);
\draw[edge, draw=kConcept!40!gmSurface] (doge) -- (ty2);
\draw[edge, draw=kConcept!40!gmSurface] (doge) -- (ty3);
\draw[edge, draw=kConcept!40!gmSurface] (doge) -- (ty4);
\draw[edge, draw=kConcept!25!gmSurface, line width=.25pt] (doge) -- (tg0);
\draw[edge, draw=kConcept!25!gmSurface, line width=.25pt] (doge) -- (tg1);
\draw[edge, draw=kConcept!25!gmSurface, line width=.25pt] (doge) -- (tg2);
\draw[edge, draw=kConcept!25!gmSurface, line width=.25pt] (doge) -- (tg3);
\draw[edge, draw=kConcept!25!gmSurface, line width=.25pt] (doge) -- (tg4);
\draw[edge, draw=kConcept!25!gmSurface, line width=.25pt] (doge) -- (tg5);
\draw[edge, draw=kConcept!25!gmSurface, line width=.25pt] (doge) -- (tg6);
\draw[edge, draw=kConcept!25!gmSurface, line width=.25pt] (doge) -- (tg7);
\draw[edge, draw=kConcept!25!gmSurface, line width=.25pt] (doge) -- (tg8);
\draw[edge, draw=kConcept!25!gmSurface, line width=.25pt] (doge) -- (tg9);
\draw[edge, draw=kConcept!25!gmSurface, line width=.25pt] (doge) -- (tg10);
\draw[edge, draw=kConcept!25!gmSurface, line width=.25pt] (doge) -- (tg11);
\draw[edge, draw=kConcept!25!gmSurface, line width=.25pt] (doge) -- (tg12);
\draw[edge, draw=kConcept!25!gmSurface, line width=.25pt] (doge) -- (tg13);
\draw[edge, draw=kConcept!25!gmSurface, line width=.25pt] (doge) -- (tg14);
\draw[edge, draw=kConcept!25!gmSurface, line width=.25pt] (doge) -- (tg15);
\draw[edge, draw=kConcept!25!gmSurface, line width=.25pt] (doge) -- (tg16);
\draw[edge, draw=kConcept!25!gmSurface, line width=.25pt] (doge) -- (tg17);
\draw[edge, draw=kConcept!25!gmSurface, line width=.25pt] (doge) -- (tg18);
\draw[edge, draw=kConcept!25!gmSurface, line width=.25pt] (doge) -- (tg19);
\draw[edge, draw=kConcept!25!gmSurface, line width=.25pt] (doge) -- (tg20);
\draw[edge, draw=kEvent!40!gmSurface, line width=.3pt] (doge) -- (ev0);
\draw[edge, draw=kEvent!40!gmSurface, line width=.3pt] (doge) -- (ev1);
\draw[draw=kEvent, line width=.5pt, -{Stealth[length=1mm]}] (ev0) -- (ev1);
\draw[edge, draw=kEvent!40!gmSurface, line width=.3pt] (doge) -- (ev2);
\draw[draw=kEvent, line width=.5pt, -{Stealth[length=1mm]}] (ev1) -- (ev2);
\draw[edge, draw=kEvent!40!gmSurface, line width=.3pt] (doge) -- (ev3);
\draw[draw=kEvent, line width=.5pt, -{Stealth[length=1mm]}] (ev2) -- (ev3);
\draw[edge, draw=kEvent!40!gmSurface, line width=.3pt] (doge) -- (ev4);
\draw[draw=kEvent, line width=.5pt, -{Stealth[length=1mm]}] (ev3) -- (ev4);
\draw[edge, draw=kEvent!40!gmSurface, line width=.3pt] (doge) -- (ev5);
\draw[draw=kEvent, line width=.5pt, -{Stealth[length=1mm]}] (ev4) -- (ev5);
\draw[edge, draw=kEvent!40!gmSurface, line width=.3pt] (doge) -- (ev6);
\draw[draw=kEvent, line width=.5pt, -{Stealth[length=1mm]}] (ev5) -- (ev6);
\draw[edge, draw=kEvent!40!gmSurface, line width=.3pt] (doge) -- (ev7);
\draw[draw=kEvent, line width=.5pt, -{Stealth[length=1mm]}] (ev6) -- (ev7);
\draw[edge, draw=kEvent!40!gmSurface, line width=.3pt] (doge) -- (ev8);
\draw[draw=kEvent, line width=.5pt, -{Stealth[length=1mm]}] (ev7) -- (ev8);
\draw[edge, draw=kEvent!40!gmSurface, line width=.3pt] (doge) -- (ev9);
\draw[draw=kEvent, line width=.5pt, -{Stealth[length=1mm]}] (ev8) -- (ev9);
\draw[edge, draw=kEvent!40!gmSurface, line width=.3pt] (doge) -- (ev10);
\draw[draw=kEvent, line width=.5pt, -{Stealth[length=1mm]}] (ev9) -- (ev10);
\draw[edge, draw=kEvent!40!gmSurface, line width=.3pt] (doge) -- (ev11);
\draw[draw=kEvent, line width=.5pt, -{Stealth[length=1mm]}] (ev10) -- (ev11);
\draw[edge, draw=kEvent!40!gmSurface, line width=.3pt] (doge) -- (ev12);
\draw[draw=kEvent, line width=.5pt, -{Stealth[length=1mm]}] (ev11) -- (ev12);
\draw[edge, draw=kEvent!40!gmSurface, line width=.3pt] (doge) -- (ev13);
\draw[draw=kEvent, line width=.5pt, -{Stealth[length=1mm]}] (ev12) -- (ev13);
\draw[edge, draw=kEvent!40!gmSurface, line width=.3pt] (doge) -- (ev14);
\draw[draw=kEvent, line width=.5pt, -{Stealth[length=1mm]}] (ev13) -- (ev14);
\draw[edge, draw=kEvent!40!gmSurface, line width=.3pt] (doge) -- (ev15);
\draw[draw=kEvent, line width=.5pt, -{Stealth[length=1mm]}] (ev14) -- (ev15);
\end{pgfonlayer}
\end{tikzpicture}
}\par\vskip1mm
  % generated by scripts/figures.py — do not edit
\begin{tikzpicture}[x=1cm]
\fill[kFrame] (0.00,0) circle (.75mm); \node[anchor=west, font=\scriptsize, text=gmInk2] at (0.10,0) {frames};
\fill[kWikidata] (2.10,0) circle (.75mm); \node[anchor=west, font=\scriptsize, text=gmInk2] at (2.20,0) {Wikidata items};
\fill[kTemplate] (4.20,0) circle (.75mm); \node[anchor=west, font=\scriptsize, text=gmInk2] at (4.30,0) {templates};
\fill[kEvent] (6.30,0) circle (.75mm); \node[anchor=west, font=\scriptsize, text=gmInk2] at (6.40,0) {events};
\fill[kConcept] (8.40,0) circle (.75mm); \node[anchor=west, font=\scriptsize, text=gmInk2] at (8.50,0) {types and tags};
\end{tikzpicture}

  \note{
  \begin{itemize}
    \item Layman version: every fact from the previous chapters is now an arrow from Doge.
    \item Up: the series it belongs to --- Interior Monologue Captioning, Image Macros, Memes.
      Upper left: the memes that grew out of it (Ironic Doge Memes, Shiba Inus, Doge 2, the
      death of Kabosu). Left, dashed: its siblings in the same series.
    \item Right: the Wikidata items its text names. Lower right: its templates, and what their
      images show --- some of those are the same items the text names (Shiba Inu, Doge): text
      and picture agree, and the graph shows it.
    \item Bottom: its 16 events, one chain. Lower left: its types and 22 tags.
    \item This is the first query of KG\_QUERIES (``one frame as a graph''), drawn from the live graph.
  \end{itemize}}
\end{frame}

\begin{frame}{And that was only part of it}
  \centering
  \resizebox{!}{6.2cm}{% generated by scripts/figures.py — do not edit
\begin{tikzpicture}[x=1cm, y=1cm]
\draw[gmBase!60!gmSurface, line width=.15pt] (0,0) -- (0.00,1.83);
\draw[gmBase!60!gmSurface, line width=.15pt] (0,0) -- (0.04,2.08);
\draw[gmBase!60!gmSurface, line width=.15pt] (0,0) -- (0.10,2.33);
\draw[gmBase!60!gmSurface, line width=.15pt] (0,0) -- (0.12,1.83);
\draw[gmBase!60!gmSurface, line width=.15pt] (0,0) -- (0.18,2.08);
\draw[gmBase!60!gmSurface, line width=.15pt] (0,0) -- (0.25,2.32);
\draw[gmBase!60!gmSurface, line width=.15pt] (0,0) -- (0.23,1.82);
\draw[gmBase!60!gmSurface, line width=.15pt] (0,0) -- (0.31,2.07);
\draw[gmBase!60!gmSurface, line width=.15pt] (0,0) -- (0.40,2.31);
\draw[gmBase!60!gmSurface, line width=.15pt] (0,0) -- (0.35,1.81);
\draw[gmBase!60!gmSurface, line width=.15pt] (0,0) -- (0.44,2.05);
\draw[gmBase!60!gmSurface, line width=.15pt] (0,0) -- (0.54,2.29);
\draw[gmBase!60!gmSurface, line width=.15pt] (0,0) -- (0.46,1.80);
\draw[gmBase!60!gmSurface, line width=.15pt] (0,0) -- (0.57,2.03);
\draw[gmBase!60!gmSurface, line width=.15pt] (0,0) -- (0.69,2.27);
\draw[gmBase!60!gmSurface, line width=.15pt] (0,0) -- (0.58,1.78);
\draw[gmBase!60!gmSurface, line width=.15pt] (0,0) -- (0.70,2.01);
\draw[gmBase!60!gmSurface, line width=.15pt] (0,0) -- (0.83,2.24);
\draw[gmBase!60!gmSurface, line width=.15pt] (0,0) -- (0.69,1.75);
\draw[gmBase!60!gmSurface, line width=.15pt] (0,0) -- (0.83,1.98);
\draw[gmBase!60!gmSurface, line width=.15pt] (0,0) -- (0.97,2.21);
\draw[gmBase!60!gmSurface, line width=.15pt] (0,0) -- (0.80,1.72);
\draw[gmBase!60!gmSurface, line width=.15pt] (0,0) -- (0.95,1.95);
\draw[gmBase!60!gmSurface, line width=.15pt] (0,0) -- (1.11,2.16);
\draw[gmBase!60!gmSurface, line width=.15pt] (0,0) -- (0.91,1.69);
\draw[gmBase!60!gmSurface, line width=.15pt] (0,0) -- (1.08,1.91);
\draw[gmBase!60!gmSurface, line width=.15pt] (0,0) -- (1.25,2.12);
\draw[gmBase!60!gmSurface, line width=.15pt] (0,0) -- (1.02,1.65);
\draw[gmBase!60!gmSurface, line width=.15pt] (0,0) -- (1.20,1.86);
\draw[gmBase!60!gmSurface, line width=.15pt] (0,0) -- (1.38,2.07);
\draw[gmBase!60!gmSurface, line width=.15pt] (0,0) -- (1.12,1.61);
\draw[gmBase!60!gmSurface, line width=.15pt] (0,0) -- (1.31,1.81);
\draw[gmBase!60!gmSurface, line width=.15pt] (0,0) -- (1.51,2.01);
\draw[gmBase!60!gmSurface, line width=.15pt] (0,0) -- (1.22,1.57);
\draw[gmBase!60!gmSurface, line width=.15pt] (0,0) -- (1.43,1.76);
\draw[gmBase!60!gmSurface, line width=.15pt] (0,0) -- (1.64,1.95);
\draw[gmBase!60!gmSurface, line width=.15pt] (0,0) -- (1.32,1.52);
\draw[gmBase!60!gmSurface, line width=.15pt] (0,0) -- (1.54,1.70);
\draw[gmBase!60!gmSurface, line width=.15pt] (0,0) -- (1.76,1.88);
\draw[gmBase!60!gmSurface, line width=.15pt] (0,0) -- (1.42,1.46);
\draw[gmBase!60!gmSurface, line width=.15pt] (0,0) -- (1.64,1.64);
\draw[gmBase!60!gmSurface, line width=.15pt] (0,0) -- (1.88,1.81);
\draw[gmBase!60!gmSurface, line width=.15pt] (0,0) -- (1.51,1.41);
\draw[gmBase!60!gmSurface, line width=.15pt] (0,0) -- (1.75,1.58);
\draw[gmBase!60!gmSurface, line width=.15pt] (0,0) -- (1.99,1.74);
\draw[gmBase!60!gmSurface, line width=.15pt] (0,0) -- (1.60,1.35);
\draw[gmBase!60!gmSurface, line width=.15pt] (0,0) -- (1.84,1.51);
\draw[gmBase!60!gmSurface, line width=.15pt] (0,0) -- (2.10,1.66);
\draw[gmBase!60!gmSurface, line width=.15pt] (0,0) -- (1.68,1.28);
\draw[gmBase!60!gmSurface, line width=.15pt] (0,0) -- (1.94,1.43);
\draw[gmBase!60!gmSurface, line width=.15pt] (0,0) -- (2.20,1.58);
\draw[gmBase!60!gmSurface, line width=.15pt] (0,0) -- (1.76,1.22);
\draw[gmBase!60!gmSurface, line width=.15pt] (0,0) -- (2.03,1.36);
\draw[gmBase!60!gmSurface, line width=.15pt] (0,0) -- (2.30,1.49);
\draw[gmBase!60!gmSurface, line width=.15pt] (0,0) -- (1.83,1.15);
\draw[gmBase!60!gmSurface, line width=.15pt] (0,0) -- (2.11,1.28);
\draw[gmBase!60!gmSurface, line width=.15pt] (0,0) -- (2.39,1.40);
\draw[gmBase!60!gmSurface, line width=.15pt] (0,0) -- (1.90,1.07);
\draw[gmBase!60!gmSurface, line width=.15pt] (0,0) -- (2.19,1.19);
\draw[gmBase!60!gmSurface, line width=.15pt] (0,0) -- (2.48,1.30);
\draw[gmBase!60!gmSurface, line width=.15pt] (0,0) -- (1.97,1.00);
\draw[gmBase!60!gmSurface, line width=.15pt] (0,0) -- (2.26,1.11);
\draw[gmBase!60!gmSurface, line width=.15pt] (0,0) -- (2.56,1.20);
\draw[gmBase!60!gmSurface, line width=.15pt] (0,0) -- (2.03,0.92);
\draw[gmBase!60!gmSurface, line width=.15pt] (0,0) -- (2.33,1.02);
\draw[gmBase!60!gmSurface, line width=.15pt] (0,0) -- (2.63,1.10);
\draw[gmBase!60!gmSurface, line width=.15pt] (0,0) -- (2.09,0.84);
\draw[gmBase!60!gmSurface, line width=.15pt] (0,0) -- (2.39,0.92);
\draw[gmBase!60!gmSurface, line width=.15pt] (0,0) -- (2.70,1.00);
\draw[gmBase!60!gmSurface, line width=.15pt] (0,0) -- (2.14,0.76);
\draw[gmBase!60!gmSurface, line width=.15pt] (0,0) -- (2.45,0.83);
\draw[gmBase!60!gmSurface, line width=.15pt] (0,0) -- (2.76,0.89);
\draw[gmBase!60!gmSurface, line width=.15pt] (0,0) -- (2.18,0.68);
\draw[gmBase!60!gmSurface, line width=.15pt] (0,0) -- (2.50,0.73);
\draw[gmBase!60!gmSurface, line width=.15pt] (0,0) -- (2.81,0.79);
\draw[gmBase!60!gmSurface, line width=.15pt] (0,0) -- (2.23,0.59);
\draw[gmBase!60!gmSurface, line width=.15pt] (0,0) -- (2.54,0.64);
\draw[gmBase!60!gmSurface, line width=.15pt] (0,0) -- (2.86,0.68);
\draw[gmBase!60!gmSurface, line width=.15pt] (0,0) -- (2.26,0.50);
\draw[gmBase!60!gmSurface, line width=.15pt] (0,0) -- (2.58,0.54);
\draw[gmBase!60!gmSurface, line width=.15pt] (0,0) -- (2.90,0.56);
\draw[gmBase!60!gmSurface, line width=.15pt] (0,0) -- (2.29,0.41);
\draw[gmBase!60!gmSurface, line width=.15pt] (0,0) -- (2.61,0.44);
\draw[gmBase!60!gmSurface, line width=.15pt] (0,0) -- (2.93,0.45);
\draw[gmBase!60!gmSurface, line width=.15pt] (0,0) -- (2.31,0.32);
\draw[gmBase!60!gmSurface, line width=.15pt] (0,0) -- (2.64,0.34);
\draw[gmBase!60!gmSurface, line width=.15pt] (0,0) -- (2.96,0.34);
\draw[gmBase!60!gmSurface, line width=.15pt] (0,0) -- (2.33,0.23);
\draw[gmBase!60!gmSurface, line width=.15pt] (0,0) -- (2.65,0.23);
\draw[gmBase!60!gmSurface, line width=.15pt] (0,0) -- (2.98,0.22);
\draw[gmBase!60!gmSurface, line width=.15pt] (0,0) -- (2.34,0.14);
\draw[gmBase!60!gmSurface, line width=.15pt] (0,0) -- (2.66,0.13);
\draw[gmBase!60!gmSurface, line width=.15pt] (0,0) -- (2.99,0.11);
\draw[gmBase!60!gmSurface, line width=.15pt] (0,0) -- (2.35,0.05);
\draw[gmBase!60!gmSurface, line width=.15pt] (0,0) -- (2.67,0.03);
\draw[gmBase!60!gmSurface, line width=.15pt] (0,0) -- (2.99,-0.01);
\draw[gmBase!60!gmSurface, line width=.15pt] (0,0) -- (2.35,-0.04);
\draw[gmBase!60!gmSurface, line width=.15pt] (0,0) -- (2.67,-0.08);
\draw[gmBase!60!gmSurface, line width=.15pt] (0,0) -- (2.99,-0.13);
\draw[gmBase!60!gmSurface, line width=.15pt] (0,0) -- (2.34,-0.13);
\draw[gmBase!60!gmSurface, line width=.15pt] (0,0) -- (2.66,-0.18);
\draw[gmBase!60!gmSurface, line width=.15pt] (0,0) -- (2.97,-0.24);
\draw[gmBase!60!gmSurface, line width=.15pt] (0,0) -- (2.33,-0.22);
\draw[gmBase!60!gmSurface, line width=.15pt] (0,0) -- (2.65,-0.28);
\draw[gmBase!60!gmSurface, line width=.15pt] (0,0) -- (2.95,-0.36);
\draw[gmBase!60!gmSurface, line width=.15pt] (0,0) -- (2.32,-0.31);
\draw[gmBase!60!gmSurface, line width=.15pt] (0,0) -- (2.62,-0.39);
\draw[gmBase!60!gmSurface, line width=.15pt] (0,0) -- (2.93,-0.47);
\draw[gmBase!60!gmSurface, line width=.15pt] (0,0) -- (2.29,-0.40);
\draw[gmBase!60!gmSurface, line width=.15pt] (0,0) -- (2.60,-0.49);
\draw[gmBase!60!gmSurface, line width=.15pt] (0,0) -- (2.89,-0.58);
\draw[gmBase!60!gmSurface, line width=.15pt] (0,0) -- (2.27,-0.49);
\draw[gmBase!60!gmSurface, line width=.15pt] (0,0) -- (2.56,-0.59);
\draw[gmBase!60!gmSurface, line width=.15pt] (0,0) -- (2.85,-0.69);
\draw[gmBase!60!gmSurface, line width=.15pt] (0,0) -- (2.23,-0.58);
\draw[gmBase!60!gmSurface, line width=.15pt] (0,0) -- (2.52,-0.69);
\draw[gmBase!60!gmSurface, line width=.15pt] (0,0) -- (2.81,-0.80);
\draw[gmBase!60!gmSurface, line width=.15pt] (0,0) -- (2.19,-0.66);
\draw[gmBase!60!gmSurface, line width=.15pt] (0,0) -- (2.47,-0.78);
\draw[gmBase!60!gmSurface, line width=.15pt] (0,0) -- (2.75,-0.91);
\draw[gmBase!60!gmSurface, line width=.15pt] (0,0) -- (2.15,-0.75);
\draw[gmBase!60!gmSurface, line width=.15pt] (0,0) -- (2.42,-0.88);
\draw[gmBase!60!gmSurface, line width=.15pt] (0,0) -- (2.69,-1.02);
\draw[gmBase!60!gmSurface, line width=.15pt] (0,0) -- (2.10,-0.83);
\draw[gmBase!60!gmSurface, line width=.15pt] (0,0) -- (2.36,-0.97);
\draw[gmBase!60!gmSurface, line width=.15pt] (0,0) -- (2.62,-1.12);
\draw[gmBase!60!gmSurface, line width=.15pt] (0,0) -- (2.04,-0.91);
\draw[gmBase!60!gmSurface, line width=.15pt] (0,0) -- (2.30,-1.06);
\draw[gmBase!60!gmSurface, line width=.15pt] (0,0) -- (2.55,-1.22);
\draw[gmBase!60!gmSurface, line width=.15pt] (0,0) -- (1.98,-0.99);
\draw[gmBase!60!gmSurface, line width=.15pt] (0,0) -- (2.23,-1.15);
\draw[gmBase!60!gmSurface, line width=.15pt] (0,0) -- (2.47,-1.32);
\draw[gmBase!60!gmSurface, line width=.15pt] (0,0) -- (1.92,-1.06);
\draw[gmBase!60!gmSurface, line width=.15pt] (0,0) -- (2.15,-1.23);
\draw[gmBase!60!gmSurface, line width=.15pt] (0,0) -- (2.38,-1.41);
\draw[gmBase!60!gmSurface, line width=.15pt] (0,0) -- (1.85,-1.13);
\draw[gmBase!60!gmSurface, line width=.15pt] (0,0) -- (2.07,-1.32);
\draw[gmBase!60!gmSurface, line width=.15pt] (0,0) -- (2.29,-1.50);
\draw[gmBase!60!gmSurface, line width=.15pt] (0,0) -- (1.77,-1.20);
\draw[gmBase!60!gmSurface, line width=.15pt] (0,0) -- (1.98,-1.39);
\draw[gmBase!60!gmSurface, line width=.15pt] (0,0) -- (2.19,-1.59);
\draw[gmBase!60!gmSurface, line width=.15pt] (0,0) -- (1.69,-1.27);
\draw[gmBase!60!gmSurface, line width=.15pt] (0,0) -- (1.89,-1.47);
\draw[gmBase!60!gmSurface, line width=.15pt] (0,0) -- (2.08,-1.67);
\draw[gmBase!60!gmSurface, line width=.15pt] (0,0) -- (1.61,-1.34);
\draw[gmBase!60!gmSurface, line width=.15pt] (0,0) -- (1.80,-1.54);
\draw[gmBase!60!gmSurface, line width=.15pt] (0,0) -- (1.97,-1.75);
\draw[gmBase!60!gmSurface, line width=.15pt] (0,0) -- (1.52,-1.40);
\draw[gmBase!60!gmSurface, line width=.15pt] (0,0) -- (1.70,-1.61);
\draw[gmBase!60!gmSurface, line width=.15pt] (0,0) -- (1.86,-1.83);
\draw[gmBase!60!gmSurface, line width=.15pt] (0,0) -- (1.43,-1.45);
\draw[gmBase!60!gmSurface, line width=.15pt] (0,0) -- (1.59,-1.67);
\draw[gmBase!60!gmSurface, line width=.15pt] (0,0) -- (1.74,-1.90);
\draw[gmBase!60!gmSurface, line width=.15pt] (0,0) -- (1.34,-1.51);
\draw[gmBase!60!gmSurface, line width=.15pt] (0,0) -- (1.48,-1.73);
\draw[gmBase!60!gmSurface, line width=.15pt] (0,0) -- (1.62,-1.96);
\draw[gmBase!60!gmSurface, line width=.15pt] (0,0) -- (1.24,-1.56);
\draw[gmBase!60!gmSurface, line width=.15pt] (0,0) -- (1.37,-1.79);
\draw[gmBase!60!gmSurface, line width=.15pt] (0,0) -- (1.49,-2.02);
\draw[gmBase!60!gmSurface, line width=.15pt] (0,0) -- (1.14,-1.60);
\draw[gmBase!60!gmSurface, line width=.15pt] (0,0) -- (1.25,-1.84);
\draw[gmBase!60!gmSurface, line width=.15pt] (0,0) -- (1.36,-2.08);
\draw[gmBase!60!gmSurface, line width=.15pt] (0,0) -- (1.03,-1.65);
\draw[gmBase!60!gmSurface, line width=.15pt] (0,0) -- (1.14,-1.88);
\draw[gmBase!60!gmSurface, line width=.15pt] (0,0) -- (1.23,-2.13);
\draw[gmBase!60!gmSurface, line width=.15pt] (0,0) -- (0.93,-1.68);
\draw[gmBase!60!gmSurface, line width=.15pt] (0,0) -- (1.01,-1.93);
\draw[gmBase!60!gmSurface, line width=.15pt] (0,0) -- (1.09,-2.17);
\draw[gmBase!60!gmSurface, line width=.15pt] (0,0) -- (0.82,-1.72);
\draw[gmBase!60!gmSurface, line width=.15pt] (0,0) -- (0.89,-1.96);
\draw[gmBase!60!gmSurface, line width=.15pt] (0,0) -- (0.95,-2.21);
\draw[gmBase!60!gmSurface, line width=.15pt] (0,0) -- (0.71,-1.75);
\draw[gmBase!60!gmSurface, line width=.15pt] (0,0) -- (0.76,-2.00);
\draw[gmBase!60!gmSurface, line width=.15pt] (0,0) -- (0.81,-2.25);
\draw[gmBase!60!gmSurface, line width=.15pt] (0,0) -- (0.60,-1.77);
\draw[gmBase!60!gmSurface, line width=.15pt] (0,0) -- (0.64,-2.02);
\draw[gmBase!60!gmSurface, line width=.15pt] (0,0) -- (0.66,-2.27);
\draw[gmBase!60!gmSurface, line width=.15pt] (0,0) -- (0.48,-1.79);
\draw[gmBase!60!gmSurface, line width=.15pt] (0,0) -- (0.51,-2.04);
\draw[gmBase!60!gmSurface, line width=.15pt] (0,0) -- (0.52,-2.30);
\draw[gmBase!60!gmSurface, line width=.15pt] (0,0) -- (0.37,-1.81);
\draw[gmBase!60!gmSurface, line width=.15pt] (0,0) -- (0.38,-2.06);
\draw[gmBase!60!gmSurface, line width=.15pt] (0,0) -- (0.37,-2.31);
\draw[gmBase!60!gmSurface, line width=.15pt] (0,0) -- (0.25,-1.82);
\draw[gmBase!60!gmSurface, line width=.15pt] (0,0) -- (0.24,-2.07);
\draw[gmBase!60!gmSurface, line width=.15pt] (0,0) -- (0.22,-2.33);
\draw[gmBase!60!gmSurface, line width=.15pt] (0,0) -- (0.14,-1.83);
\draw[gmBase!60!gmSurface, line width=.15pt] (0,0) -- (0.11,-2.08);
\draw[gmBase!60!gmSurface, line width=.15pt] (0,0) -- (0.07,-2.33);
\draw[gmBase!60!gmSurface, line width=.15pt] (0,0) -- (0.02,-1.83);
\draw[gmBase!60!gmSurface, line width=.15pt] (0,0) -- (-0.02,-2.08);
\draw[gmBase!60!gmSurface, line width=.15pt] (0,0) -- (-0.07,-2.33);
\draw[gmBase!60!gmSurface, line width=.15pt] (0,0) -- (-0.10,-1.83);
\draw[gmBase!60!gmSurface, line width=.15pt] (0,0) -- (-0.15,-2.08);
\draw[gmBase!60!gmSurface, line width=.15pt] (0,0) -- (-0.22,-2.33);
\draw[gmBase!60!gmSurface, line width=.15pt] (0,0) -- (-0.21,-1.83);
\draw[gmBase!60!gmSurface, line width=.15pt] (0,0) -- (-0.29,-2.07);
\draw[gmBase!60!gmSurface, line width=.15pt] (0,0) -- (-0.37,-2.31);
\draw[gmBase!60!gmSurface, line width=.15pt] (0,0) -- (-0.33,-1.81);
\draw[gmBase!60!gmSurface, line width=.15pt] (0,0) -- (-0.42,-2.06);
\draw[gmBase!60!gmSurface, line width=.15pt] (0,0) -- (-0.52,-2.30);
\draw[gmBase!60!gmSurface, line width=.15pt] (0,0) -- (-0.45,-1.80);
\draw[gmBase!60!gmSurface, line width=.15pt] (0,0) -- (-0.55,-2.04);
\draw[gmBase!60!gmSurface, line width=.15pt] (0,0) -- (-0.66,-2.27);
\draw[gmBase!60!gmSurface, line width=.15pt] (0,0) -- (-0.56,-1.78);
\draw[gmBase!60!gmSurface, line width=.15pt] (0,0) -- (-0.68,-2.01);
\draw[gmBase!60!gmSurface, line width=.15pt] (0,0) -- (-0.81,-2.25);
\draw[gmBase!60!gmSurface, line width=.15pt] (0,0) -- (-0.67,-1.76);
\draw[gmBase!60!gmSurface, line width=.15pt] (0,0) -- (-0.81,-1.99);
\draw[gmBase!60!gmSurface, line width=.15pt] (0,0) -- (-0.95,-2.21);
\draw[gmBase!60!gmSurface, line width=.15pt] (0,0) -- (-0.78,-1.73);
\draw[gmBase!60!gmSurface, line width=.15pt] (0,0) -- (-0.93,-1.95);
\draw[gmBase!60!gmSurface, line width=.15pt] (0,0) -- (-1.09,-2.17);
\draw[gmBase!60!gmSurface, line width=.15pt] (0,0) -- (-0.89,-1.70);
\draw[gmBase!60!gmSurface, line width=.15pt] (0,0) -- (-1.05,-1.91);
\draw[gmBase!60!gmSurface, line width=.15pt] (0,0) -- (-1.23,-2.13);
\draw[gmBase!60!gmSurface, line width=.15pt] (0,0) -- (-1.00,-1.66);
\draw[gmBase!60!gmSurface, line width=.15pt] (0,0) -- (-1.18,-1.87);
\draw[gmBase!60!gmSurface, line width=.15pt] (0,0) -- (-1.36,-2.08);
\draw[gmBase!60!gmSurface, line width=.15pt] (0,0) -- (-1.10,-1.62);
\draw[gmBase!60!gmSurface, line width=.15pt] (0,0) -- (-1.29,-1.82);
\draw[gmBase!60!gmSurface, line width=.15pt] (0,0) -- (-1.49,-2.02);
\draw[gmBase!60!gmSurface, line width=.15pt] (0,0) -- (-1.21,-1.57);
\draw[gmBase!60!gmSurface, line width=.15pt] (0,0) -- (-1.41,-1.77);
\draw[gmBase!60!gmSurface, line width=.15pt] (0,0) -- (-1.62,-1.96);
\draw[gmBase!60!gmSurface, line width=.15pt] (0,0) -- (-1.30,-1.52);
\draw[gmBase!60!gmSurface, line width=.15pt] (0,0) -- (-1.52,-1.71);
\draw[gmBase!60!gmSurface, line width=.15pt] (0,0) -- (-1.74,-1.90);
\draw[gmBase!60!gmSurface, line width=.15pt] (0,0) -- (-1.40,-1.47);
\draw[gmBase!60!gmSurface, line width=.15pt] (0,0) -- (-1.63,-1.65);
\draw[gmBase!60!gmSurface, line width=.15pt] (0,0) -- (-1.86,-1.83);
\draw[gmBase!60!gmSurface, line width=.15pt] (0,0) -- (-1.49,-1.42);
\draw[gmBase!60!gmSurface, line width=.15pt] (0,0) -- (-1.73,-1.59);
\draw[gmBase!60!gmSurface, line width=.15pt] (0,0) -- (-1.97,-1.75);
\draw[gmBase!60!gmSurface, line width=.15pt] (0,0) -- (-1.58,-1.36);
\draw[gmBase!60!gmSurface, line width=.15pt] (0,0) -- (-1.83,-1.52);
\draw[gmBase!60!gmSurface, line width=.15pt] (0,0) -- (-2.08,-1.67);
\draw[gmBase!60!gmSurface, line width=.15pt] (0,0) -- (-1.67,-1.29);
\draw[gmBase!60!gmSurface, line width=.15pt] (0,0) -- (-1.92,-1.44);
\draw[gmBase!60!gmSurface, line width=.15pt] (0,0) -- (-2.19,-1.59);
\draw[gmBase!60!gmSurface, line width=.15pt] (0,0) -- (-1.75,-1.23);
\draw[gmBase!60!gmSurface, line width=.15pt] (0,0) -- (-2.01,-1.37);
\draw[gmBase!60!gmSurface, line width=.15pt] (0,0) -- (-2.29,-1.50);
\draw[gmBase!60!gmSurface, line width=.15pt] (0,0) -- (-1.82,-1.16);
\draw[gmBase!60!gmSurface, line width=.15pt] (0,0) -- (-2.10,-1.29);
\draw[gmBase!60!gmSurface, line width=.15pt] (0,0) -- (-2.38,-1.41);
\draw[gmBase!60!gmSurface, line width=.15pt] (0,0) -- (-1.89,-1.09);
\draw[gmBase!60!gmSurface, line width=.15pt] (0,0) -- (-2.18,-1.21);
\draw[gmBase!60!gmSurface, line width=.15pt] (0,0) -- (-2.47,-1.32);
\draw[gmBase!60!gmSurface, line width=.15pt] (0,0) -- (-1.96,-1.01);
\draw[gmBase!60!gmSurface, line width=.15pt] (0,0) -- (-2.25,-1.12);
\draw[gmBase!60!gmSurface, line width=.15pt] (0,0) -- (-2.55,-1.22);
\draw[gmBase!60!gmSurface, line width=.15pt] (0,0) -- (-2.02,-0.93);
\draw[gmBase!60!gmSurface, line width=.15pt] (0,0) -- (-2.32,-1.03);
\draw[gmBase!60!gmSurface, line width=.15pt] (0,0) -- (-2.62,-1.12);
\draw[gmBase!60!gmSurface, line width=.15pt] (0,0) -- (-2.08,-0.85);
\draw[gmBase!60!gmSurface, line width=.15pt] (0,0) -- (-2.38,-0.94);
\draw[gmBase!60!gmSurface, line width=.15pt] (0,0) -- (-2.69,-1.02);
\draw[gmBase!60!gmSurface, line width=.15pt] (0,0) -- (-2.13,-0.77);
\draw[gmBase!60!gmSurface, line width=.15pt] (0,0) -- (-2.44,-0.85);
\draw[gmBase!60!gmSurface, line width=.15pt] (0,0) -- (-2.75,-0.91);
\draw[gmBase!60!gmSurface, line width=.15pt] (0,0) -- (-2.18,-0.69);
\draw[gmBase!60!gmSurface, line width=.15pt] (0,0) -- (-2.49,-0.75);
\draw[gmBase!60!gmSurface, line width=.15pt] (0,0) -- (-2.81,-0.80);
\draw[gmBase!60!gmSurface, line width=.15pt] (0,0) -- (-2.22,-0.60);
\draw[gmBase!60!gmSurface, line width=.15pt] (0,0) -- (-2.54,-0.65);
\draw[gmBase!60!gmSurface, line width=.15pt] (0,0) -- (-2.85,-0.69);
\draw[gmBase!60!gmSurface, line width=.15pt] (0,0) -- (-2.25,-0.52);
\draw[gmBase!60!gmSurface, line width=.15pt] (0,0) -- (-2.57,-0.55);
\draw[gmBase!60!gmSurface, line width=.15pt] (0,0) -- (-2.89,-0.58);
\draw[gmBase!60!gmSurface, line width=.15pt] (0,0) -- (-2.28,-0.43);
\draw[gmBase!60!gmSurface, line width=.15pt] (0,0) -- (-2.61,-0.45);
\draw[gmBase!60!gmSurface, line width=.15pt] (0,0) -- (-2.93,-0.47);
\draw[gmBase!60!gmSurface, line width=.15pt] (0,0) -- (-2.31,-0.34);
\draw[gmBase!60!gmSurface, line width=.15pt] (0,0) -- (-2.63,-0.35);
\draw[gmBase!60!gmSurface, line width=.15pt] (0,0) -- (-2.95,-0.36);
\draw[gmBase!60!gmSurface, line width=.15pt] (0,0) -- (-2.33,-0.25);
\draw[gmBase!60!gmSurface, line width=.15pt] (0,0) -- (-2.65,-0.25);
\draw[gmBase!60!gmSurface, line width=.15pt] (0,0) -- (-2.97,-0.24);
\draw[gmBase!60!gmSurface, line width=.15pt] (0,0) -- (-2.34,-0.16);
\draw[gmBase!60!gmSurface, line width=.15pt] (0,0) -- (-2.66,-0.15);
\draw[gmBase!60!gmSurface, line width=.15pt] (0,0) -- (-2.99,-0.13);
\draw[gmBase!60!gmSurface, line width=.15pt] (0,0) -- (-2.35,-0.07);
\draw[gmBase!60!gmSurface, line width=.15pt] (0,0) -- (-2.67,-0.04);
\draw[gmBase!60!gmSurface, line width=.15pt] (0,0) -- (-2.99,-0.01);
\draw[gmBase!60!gmSurface, line width=.15pt] (0,0) -- (-2.35,0.02);
\draw[gmBase!60!gmSurface, line width=.15pt] (0,0) -- (-2.67,0.06);
\draw[gmBase!60!gmSurface, line width=.15pt] (0,0) -- (-2.99,0.11);
\draw[gmBase!60!gmSurface, line width=.15pt] (0,0) -- (-2.35,0.11);
\draw[gmBase!60!gmSurface, line width=.15pt] (0,0) -- (-2.66,0.16);
\draw[gmBase!60!gmSurface, line width=.15pt] (0,0) -- (-2.98,0.22);
\draw[gmBase!60!gmSurface, line width=.15pt] (0,0) -- (-2.34,0.20);
\draw[gmBase!60!gmSurface, line width=.15pt] (0,0) -- (-2.65,0.27);
\draw[gmBase!60!gmSurface, line width=.15pt] (0,0) -- (-2.96,0.34);
\draw[gmBase!60!gmSurface, line width=.15pt] (0,0) -- (-2.32,0.29);
\draw[gmBase!60!gmSurface, line width=.15pt] (0,0) -- (-2.63,0.37);
\draw[gmBase!60!gmSurface, line width=.15pt] (0,0) -- (-2.93,0.45);
\draw[gmBase!60!gmSurface, line width=.15pt] (0,0) -- (-2.30,0.38);
\draw[gmBase!60!gmSurface, line width=.15pt] (0,0) -- (-2.60,0.47);
\draw[gmBase!60!gmSurface, line width=.15pt] (0,0) -- (-2.90,0.56);
\draw[gmBase!60!gmSurface, line width=.15pt] (0,0) -- (-2.27,0.47);
\draw[gmBase!60!gmSurface, line width=.15pt] (0,0) -- (-2.57,0.57);
\draw[gmBase!60!gmSurface, line width=.15pt] (0,0) -- (-2.86,0.68);
\draw[gmBase!60!gmSurface, line width=.15pt] (0,0) -- (-2.24,0.56);
\draw[gmBase!60!gmSurface, line width=.15pt] (0,0) -- (-2.53,0.67);
\draw[gmBase!60!gmSurface, line width=.15pt] (0,0) -- (-2.81,0.79);
\draw[gmBase!60!gmSurface, line width=.15pt] (0,0) -- (-2.20,0.65);
\draw[gmBase!60!gmSurface, line width=.15pt] (0,0) -- (-2.48,0.77);
\draw[gmBase!60!gmSurface, line width=.15pt] (0,0) -- (-2.76,0.89);
\draw[gmBase!60!gmSurface, line width=.15pt] (0,0) -- (-2.15,0.73);
\draw[gmBase!60!gmSurface, line width=.15pt] (0,0) -- (-2.43,0.86);
\draw[gmBase!60!gmSurface, line width=.15pt] (0,0) -- (-2.70,1.00);
\draw[gmBase!60!gmSurface, line width=.15pt] (0,0) -- (-2.11,0.81);
\draw[gmBase!60!gmSurface, line width=.15pt] (0,0) -- (-2.37,0.96);
\draw[gmBase!60!gmSurface, line width=.15pt] (0,0) -- (-2.63,1.10);
\draw[gmBase!60!gmSurface, line width=.15pt] (0,0) -- (-2.05,0.89);
\draw[gmBase!60!gmSurface, line width=.15pt] (0,0) -- (-2.31,1.05);
\draw[gmBase!60!gmSurface, line width=.15pt] (0,0) -- (-2.56,1.20);
\draw[gmBase!60!gmSurface, line width=.15pt] (0,0) -- (-1.99,0.97);
\draw[gmBase!60!gmSurface, line width=.15pt] (0,0) -- (-2.24,1.13);
\draw[gmBase!60!gmSurface, line width=.15pt] (0,0) -- (-2.48,1.30);
\draw[gmBase!60!gmSurface, line width=.15pt] (0,0) -- (-1.93,1.05);
\draw[gmBase!60!gmSurface, line width=.15pt] (0,0) -- (-2.16,1.22);
\draw[gmBase!60!gmSurface, line width=.15pt] (0,0) -- (-2.39,1.40);
\draw[gmBase!60!gmSurface, line width=.15pt] (0,0) -- (-1.86,1.12);
\draw[gmBase!60!gmSurface, line width=.15pt] (0,0) -- (-2.08,1.30);
\draw[gmBase!60!gmSurface, line width=.15pt] (0,0) -- (-2.30,1.49);
\draw[gmBase!60!gmSurface, line width=.15pt] (0,0) -- (-1.78,1.19);
\draw[gmBase!60!gmSurface, line width=.15pt] (0,0) -- (-2.00,1.38);
\draw[gmBase!60!gmSurface, line width=.15pt] (0,0) -- (-2.20,1.58);
\draw[gmBase!60!gmSurface, line width=.15pt] (0,0) -- (-1.71,1.26);
\draw[gmBase!60!gmSurface, line width=.15pt] (0,0) -- (-1.91,1.46);
\draw[gmBase!60!gmSurface, line width=.15pt] (0,0) -- (-2.10,1.66);
\draw[gmBase!60!gmSurface, line width=.15pt] (0,0) -- (-1.62,1.33);
\draw[gmBase!60!gmSurface, line width=.15pt] (0,0) -- (-1.81,1.53);
\draw[gmBase!60!gmSurface, line width=.15pt] (0,0) -- (-1.99,1.74);
\draw[gmBase!60!gmSurface, line width=.15pt] (0,0) -- (-1.54,1.39);
\draw[gmBase!60!gmSurface, line width=.15pt] (0,0) -- (-1.71,1.60);
\draw[gmBase!60!gmSurface, line width=.15pt] (0,0) -- (-1.88,1.81);
\draw[gmBase!60!gmSurface, line width=.15pt] (0,0) -- (-1.45,1.44);
\draw[gmBase!60!gmSurface, line width=.15pt] (0,0) -- (-1.61,1.66);
\draw[gmBase!60!gmSurface, line width=.15pt] (0,0) -- (-1.76,1.88);
\draw[gmBase!60!gmSurface, line width=.15pt] (0,0) -- (-1.35,1.50);
\draw[gmBase!60!gmSurface, line width=.15pt] (0,0) -- (-1.50,1.72);
\draw[gmBase!60!gmSurface, line width=.15pt] (0,0) -- (-1.64,1.95);
\draw[gmBase!60!gmSurface, line width=.15pt] (0,0) -- (-1.26,1.55);
\draw[gmBase!60!gmSurface, line width=.15pt] (0,0) -- (-1.39,1.78);
\draw[gmBase!60!gmSurface, line width=.15pt] (0,0) -- (-1.51,2.01);
\draw[gmBase!60!gmSurface, line width=.15pt] (0,0) -- (-1.16,1.60);
\draw[gmBase!60!gmSurface, line width=.15pt] (0,0) -- (-1.27,1.83);
\draw[gmBase!60!gmSurface, line width=.15pt] (0,0) -- (-1.38,2.07);
\draw[gmBase!60!gmSurface, line width=.15pt] (0,0) -- (-1.05,1.64);
\draw[gmBase!60!gmSurface, line width=.15pt] (0,0) -- (-1.16,1.88);
\draw[gmBase!60!gmSurface, line width=.15pt] (0,0) -- (-1.25,2.12);
\draw[gmBase!60!gmSurface, line width=.15pt] (0,0) -- (-0.95,1.68);
\draw[gmBase!60!gmSurface, line width=.15pt] (0,0) -- (-1.03,1.92);
\draw[gmBase!60!gmSurface, line width=.15pt] (0,0) -- (-1.11,2.16);
\draw[gmBase!60!gmSurface, line width=.15pt] (0,0) -- (-0.84,1.71);
\draw[gmBase!60!gmSurface, line width=.15pt] (0,0) -- (-0.91,1.96);
\draw[gmBase!60!gmSurface, line width=.15pt] (0,0) -- (-0.97,2.21);
\draw[gmBase!60!gmSurface, line width=.15pt] (0,0) -- (-0.73,1.74);
\draw[gmBase!60!gmSurface, line width=.15pt] (0,0) -- (-0.78,1.99);
\draw[gmBase!60!gmSurface, line width=.15pt] (0,0) -- (-0.83,2.24);
\draw[gmBase!60!gmSurface, line width=.15pt] (0,0) -- (-0.62,1.77);
\draw[gmBase!60!gmSurface, line width=.15pt] (0,0) -- (-0.66,2.02);
\draw[gmBase!60!gmSurface, line width=.15pt] (0,0) -- (-0.69,2.27);
\draw[gmBase!60!gmSurface, line width=.15pt] (0,0) -- (-0.50,1.79);
\draw[gmBase!60!gmSurface, line width=.15pt] (0,0) -- (-0.53,2.04);
\draw[gmBase!60!gmSurface, line width=.15pt] (0,0) -- (-0.54,2.29);
\draw[gmBase!60!gmSurface, line width=.15pt] (0,0) -- (-0.39,1.81);
\draw[gmBase!60!gmSurface, line width=.15pt] (0,0) -- (-0.40,2.06);
\draw[gmBase!60!gmSurface, line width=.15pt] (0,0) -- (-0.40,2.31);
\draw[gmBase!60!gmSurface, line width=.15pt] (0,0) -- (-0.27,1.82);
\draw[gmBase!60!gmSurface, line width=.15pt] (0,0) -- (-0.27,2.07);
\draw[gmBase!60!gmSurface, line width=.15pt] (0,0) -- (-0.25,2.32);
\draw[gmBase!60!gmSurface, line width=.15pt] (0,0) -- (-0.16,1.83);
\draw[gmBase!60!gmSurface, line width=.15pt] (0,0) -- (-0.13,2.08);
\draw[gmBase!60!gmSurface, line width=.15pt] (0,0) -- (-0.10,2.33);
\draw[gmBase!60!gmSurface, line width=.15pt] (0,0) -- (-0.04,1.83);
\fill[kFrame] (0.00,1.83) circle (.05);
\fill[kFrame] (0.04,2.08) circle (.05);
\fill[kFrame] (0.10,2.33) circle (.05);
\fill[kFrame] (0.12,1.83) circle (.05);
\fill[kFrame] (0.18,2.08) circle (.05);
\fill[kFrame] (0.25,2.32) circle (.05);
\fill[kFrame] (0.23,1.82) circle (.05);
\fill[kFrame] (0.31,2.07) circle (.05);
\fill[kFrame] (0.40,2.31) circle (.05);
\fill[kFrame] (0.35,1.81) circle (.05);
\fill[kFrame] (0.44,2.05) circle (.05);
\fill[kFrame] (0.54,2.29) circle (.05);
\fill[kFrame] (0.46,1.80) circle (.05);
\fill[kFrame] (0.57,2.03) circle (.05);
\fill[kFrame] (0.69,2.27) circle (.05);
\fill[kFrame] (0.58,1.78) circle (.05);
\fill[kFrame] (0.70,2.01) circle (.05);
\fill[kFrame] (0.83,2.24) circle (.05);
\fill[kFrame] (0.69,1.75) circle (.05);
\fill[kFrame] (0.83,1.98) circle (.05);
\fill[kFrame] (0.97,2.21) circle (.05);
\fill[kFrame] (0.80,1.72) circle (.05);
\fill[kFrame] (0.95,1.95) circle (.05);
\fill[kFrame] (1.11,2.16) circle (.05);
\fill[kFrame] (0.91,1.69) circle (.05);
\fill[kFrame] (1.08,1.91) circle (.05);
\fill[kFrame] (1.25,2.12) circle (.05);
\fill[kFrame] (1.02,1.65) circle (.05);
\fill[kFrame] (1.20,1.86) circle (.05);
\fill[kFrame] (1.38,2.07) circle (.05);
\fill[kFrame] (1.12,1.61) circle (.05);
\fill[kFrame] (1.31,1.81) circle (.05);
\fill[kFrame] (1.51,2.01) circle (.05);
\fill[kFrame] (1.22,1.57) circle (.05);
\fill[kFrame] (1.43,1.76) circle (.05);
\fill[kFrame] (1.64,1.95) circle (.05);
\fill[kFrame] (1.32,1.52) circle (.05);
\fill[kFrame] (1.54,1.70) circle (.05);
\fill[kFrame] (1.76,1.88) circle (.05);
\fill[kFrame] (1.42,1.46) circle (.05);
\fill[kFrame] (1.64,1.64) circle (.05);
\fill[kFrame] (1.88,1.81) circle (.05);
\fill[kFrame] (1.51,1.41) circle (.05);
\fill[kFrame] (1.75,1.58) circle (.05);
\fill[kFrame] (1.99,1.74) circle (.05);
\fill[kFrame] (1.60,1.35) circle (.05);
\fill[kFrame] (1.84,1.51) circle (.05);
\fill[kFrame] (2.10,1.66) circle (.05);
\fill[kFrame] (1.68,1.28) circle (.05);
\fill[kFrame] (1.94,1.43) circle (.05);
\fill[kFrame] (2.20,1.58) circle (.05);
\fill[kFrame] (1.76,1.22) circle (.05);
\fill[kFrame] (2.03,1.36) circle (.05);
\fill[kFrame] (2.30,1.49) circle (.05);
\fill[kFrame] (1.83,1.15) circle (.05);
\fill[kFrame] (2.11,1.28) circle (.05);
\fill[kFrame] (2.39,1.40) circle (.05);
\fill[kFrame] (1.90,1.07) circle (.05);
\fill[kFrame] (2.19,1.19) circle (.05);
\fill[kFrame] (2.48,1.30) circle (.05);
\fill[kFrame] (1.97,1.00) circle (.05);
\fill[kFrame] (2.26,1.11) circle (.05);
\fill[kFrame] (2.56,1.20) circle (.05);
\fill[kFrame] (2.03,0.92) circle (.05);
\fill[kFrame] (2.33,1.02) circle (.05);
\fill[kFrame] (2.63,1.10) circle (.05);
\fill[kFrame] (2.09,0.84) circle (.05);
\fill[kFrame] (2.39,0.92) circle (.05);
\fill[kFrame] (2.70,1.00) circle (.05);
\fill[kFrame] (2.14,0.76) circle (.05);
\fill[kFrame] (2.45,0.83) circle (.05);
\fill[kFrame] (2.76,0.89) circle (.05);
\fill[kFrame] (2.18,0.68) circle (.05);
\fill[kFrame] (2.50,0.73) circle (.05);
\fill[kFrame] (2.81,0.79) circle (.05);
\fill[kFrame] (2.23,0.59) circle (.05);
\fill[kFrame] (2.54,0.64) circle (.05);
\fill[kFrame] (2.86,0.68) circle (.05);
\fill[kFrame] (2.26,0.50) circle (.05);
\fill[kFrame] (2.58,0.54) circle (.05);
\fill[kFrame] (2.90,0.56) circle (.05);
\fill[kFrame] (2.29,0.41) circle (.05);
\fill[kFrame] (2.61,0.44) circle (.05);
\fill[kFrame] (2.93,0.45) circle (.05);
\fill[kFrame] (2.31,0.32) circle (.05);
\fill[kFrame] (2.64,0.34) circle (.05);
\fill[kFrame] (2.96,0.34) circle (.05);
\fill[kFrame] (2.33,0.23) circle (.05);
\fill[kFrame] (2.65,0.23) circle (.05);
\fill[kFrame] (2.98,0.22) circle (.05);
\fill[kFrame] (2.34,0.14) circle (.05);
\fill[kFrame] (2.66,0.13) circle (.05);
\fill[kFrame] (2.99,0.11) circle (.05);
\fill[kFrame] (2.35,0.05) circle (.05);
\fill[kFrame] (2.67,0.03) circle (.05);
\fill[kFrame] (2.99,-0.01) circle (.05);
\fill[kFrame] (2.35,-0.04) circle (.05);
\fill[kFrame] (2.67,-0.08) circle (.05);
\fill[kFrame] (2.99,-0.13) circle (.05);
\fill[kFrame] (2.34,-0.13) circle (.05);
\fill[kFrame] (2.66,-0.18) circle (.05);
\fill[kFrame] (2.97,-0.24) circle (.05);
\fill[kFrame] (2.33,-0.22) circle (.05);
\fill[kFrame] (2.65,-0.28) circle (.05);
\fill[kFrame] (2.95,-0.36) circle (.05);
\fill[kFrame] (2.32,-0.31) circle (.05);
\fill[kFrame] (2.62,-0.39) circle (.05);
\fill[kFrame] (2.93,-0.47) circle (.05);
\fill[kFrame] (2.29,-0.40) circle (.05);
\fill[kFrame] (2.60,-0.49) circle (.05);
\fill[kFrame] (2.89,-0.58) circle (.05);
\fill[kFrame] (2.27,-0.49) circle (.05);
\fill[kFrame] (2.56,-0.59) circle (.05);
\fill[kFrame] (2.85,-0.69) circle (.05);
\fill[kFrame] (2.23,-0.58) circle (.05);
\fill[kFrame] (2.52,-0.69) circle (.05);
\fill[kFrame] (2.81,-0.80) circle (.05);
\fill[kFrame] (2.19,-0.66) circle (.05);
\fill[kFrame] (2.47,-0.78) circle (.05);
\fill[kFrame] (2.75,-0.91) circle (.05);
\fill[kFrame] (2.15,-0.75) circle (.05);
\fill[kFrame] (2.42,-0.88) circle (.05);
\fill[kFrame] (2.69,-1.02) circle (.05);
\fill[kFrame] (2.10,-0.83) circle (.05);
\fill[kFrame] (2.36,-0.97) circle (.05);
\fill[kFrame] (2.62,-1.12) circle (.05);
\fill[kFrame] (2.04,-0.91) circle (.05);
\fill[kFrame] (2.30,-1.06) circle (.05);
\fill[kFrame] (2.55,-1.22) circle (.05);
\fill[kFrame] (1.98,-0.99) circle (.05);
\fill[kFrame] (2.23,-1.15) circle (.05);
\fill[kFrame] (2.47,-1.32) circle (.05);
\fill[kFrame] (1.92,-1.06) circle (.05);
\fill[kFrame] (2.15,-1.23) circle (.05);
\fill[kFrame] (2.38,-1.41) circle (.05);
\fill[kFrame] (1.85,-1.13) circle (.05);
\fill[kFrame] (2.07,-1.32) circle (.05);
\fill[kFrame] (2.29,-1.50) circle (.05);
\fill[kFrame] (1.77,-1.20) circle (.05);
\fill[kFrame] (1.98,-1.39) circle (.05);
\fill[kFrame] (2.19,-1.59) circle (.05);
\fill[kFrame] (1.69,-1.27) circle (.05);
\fill[kFrame] (1.89,-1.47) circle (.05);
\fill[kFrame] (2.08,-1.67) circle (.05);
\fill[kFrame] (1.61,-1.34) circle (.05);
\fill[kFrame] (1.80,-1.54) circle (.05);
\fill[kFrame] (1.97,-1.75) circle (.05);
\fill[kFrame] (1.52,-1.40) circle (.05);
\fill[kFrame] (1.70,-1.61) circle (.05);
\fill[kFrame] (1.86,-1.83) circle (.05);
\fill[kFrame] (1.43,-1.45) circle (.05);
\fill[kFrame] (1.59,-1.67) circle (.05);
\fill[kFrame] (1.74,-1.90) circle (.05);
\fill[kFrame] (1.34,-1.51) circle (.05);
\fill[kFrame] (1.48,-1.73) circle (.05);
\fill[kFrame] (1.62,-1.96) circle (.05);
\fill[kEvent] (1.24,-1.56) circle (.05);
\fill[kEvent] (1.37,-1.79) circle (.05);
\fill[kEvent] (1.49,-2.02) circle (.05);
\fill[kEvent] (1.14,-1.60) circle (.05);
\fill[kEvent] (1.25,-1.84) circle (.05);
\fill[kEvent] (1.36,-2.08) circle (.05);
\fill[kEvent] (1.03,-1.65) circle (.05);
\fill[kEvent] (1.14,-1.88) circle (.05);
\fill[kEvent] (1.23,-2.13) circle (.05);
\fill[kEvent] (0.93,-1.68) circle (.05);
\fill[kEvent] (1.01,-1.93) circle (.05);
\fill[kEvent] (1.09,-2.17) circle (.05);
\fill[kEvent] (0.82,-1.72) circle (.05);
\fill[kEvent] (0.89,-1.96) circle (.05);
\fill[kEvent] (0.95,-2.21) circle (.05);
\fill[kEvent] (0.71,-1.75) circle (.05);
\fill[kEvent] (0.76,-2.00) circle (.05);
\fill[kEvent] (0.81,-2.25) circle (.05);
\fill[kEvent] (0.60,-1.77) circle (.05);
\fill[kEvent] (0.64,-2.02) circle (.05);
\fill[kEvent] (0.66,-2.27) circle (.05);
\fill[kEvent] (0.48,-1.79) circle (.05);
\fill[kEvent] (0.51,-2.04) circle (.05);
\fill[kEvent] (0.52,-2.30) circle (.05);
\fill[kEvent] (0.37,-1.81) circle (.05);
\fill[kEvent] (0.38,-2.06) circle (.05);
\fill[kEvent] (0.37,-2.31) circle (.05);
\fill[kEvent] (0.25,-1.82) circle (.05);
\fill[kEvent] (0.24,-2.07) circle (.05);
\fill[kEvent] (0.22,-2.33) circle (.05);
\fill[kEvent] (0.14,-1.83) circle (.05);
\fill[kEvent] (0.11,-2.08) circle (.05);
\fill[kEvent] (0.07,-2.33) circle (.05);
\fill[kEvent] (0.02,-1.83) circle (.05);
\fill[kEvent] (-0.02,-2.08) circle (.05);
\fill[kEvent] (-0.07,-2.33) circle (.05);
\fill[kEvent] (-0.10,-1.83) circle (.05);
\fill[kEvent] (-0.15,-2.08) circle (.05);
\fill[kEvent] (-0.22,-2.33) circle (.05);
\fill[kEvent] (-0.21,-1.83) circle (.05);
\fill[kEvent] (-0.29,-2.07) circle (.05);
\fill[kEvent] (-0.37,-2.31) circle (.05);
\fill[kEvent] (-0.33,-1.81) circle (.05);
\fill[kEvent] (-0.42,-2.06) circle (.05);
\fill[kEvent] (-0.52,-2.30) circle (.05);
\fill[kEvent] (-0.45,-1.80) circle (.05);
\fill[kEvent] (-0.55,-2.04) circle (.05);
\fill[kEvent] (-0.66,-2.27) circle (.05);
\fill[kEvent] (-0.56,-1.78) circle (.05);
\fill[kEvent] (-0.68,-2.01) circle (.05);
\fill[kEvent] (-0.81,-2.25) circle (.05);
\fill[kEvent] (-0.67,-1.76) circle (.05);
\fill[kEvent] (-0.81,-1.99) circle (.05);
\fill[kEvent] (-0.95,-2.21) circle (.05);
\fill[kEvent] (-0.78,-1.73) circle (.05);
\fill[kEvent] (-0.93,-1.95) circle (.05);
\fill[kEvent] (-1.09,-2.17) circle (.05);
\fill[kEvent] (-0.89,-1.70) circle (.05);
\fill[kEvent] (-1.05,-1.91) circle (.05);
\fill[kEvent] (-1.23,-2.13) circle (.05);
\fill[kEvent] (-1.00,-1.66) circle (.05);
\fill[kEvent] (-1.18,-1.87) circle (.05);
\fill[kEvent] (-1.36,-2.08) circle (.05);
\fill[kEvent] (-1.10,-1.62) circle (.05);
\fill[kDoc] (-1.29,-1.82) circle (.05);
\fill[kDoc] (-1.49,-2.02) circle (.05);
\fill[kDoc] (-1.21,-1.57) circle (.05);
\fill[kDoc] (-1.41,-1.77) circle (.05);
\fill[kDoc] (-1.62,-1.96) circle (.05);
\fill[kDoc] (-1.30,-1.52) circle (.05);
\fill[kDoc] (-1.52,-1.71) circle (.05);
\fill[kDoc] (-1.74,-1.90) circle (.05);
\fill[kDoc] (-1.40,-1.47) circle (.05);
\fill[kDoc] (-1.63,-1.65) circle (.05);
\fill[kDoc] (-1.86,-1.83) circle (.05);
\fill[kDoc] (-1.49,-1.42) circle (.05);
\fill[kDoc] (-1.73,-1.59) circle (.05);
\fill[kDoc] (-1.97,-1.75) circle (.05);
\fill[kDoc] (-1.58,-1.36) circle (.05);
\fill[kDoc] (-1.83,-1.52) circle (.05);
\fill[kDoc] (-2.08,-1.67) circle (.05);
\fill[kDoc] (-1.67,-1.29) circle (.05);
\fill[kDoc] (-1.92,-1.44) circle (.05);
\fill[kDoc] (-2.19,-1.59) circle (.05);
\fill[kDoc] (-1.75,-1.23) circle (.05);
\fill[kDoc] (-2.01,-1.37) circle (.05);
\fill[kDoc] (-2.29,-1.50) circle (.05);
\fill[kDoc] (-1.82,-1.16) circle (.05);
\fill[kDoc] (-2.10,-1.29) circle (.05);
\fill[kDoc] (-2.38,-1.41) circle (.05);
\fill[kDoc] (-1.89,-1.09) circle (.05);
\fill[kDoc] (-2.18,-1.21) circle (.05);
\fill[kDoc] (-2.47,-1.32) circle (.05);
\fill[kDoc] (-1.96,-1.01) circle (.05);
\fill[kDoc] (-2.25,-1.12) circle (.05);
\fill[kDoc] (-2.55,-1.22) circle (.05);
\fill[kDoc] (-2.02,-0.93) circle (.05);
\fill[kDoc] (-2.32,-1.03) circle (.05);
\fill[kDoc] (-2.62,-1.12) circle (.05);
\fill[kDoc] (-2.08,-0.85) circle (.05);
\fill[kDoc] (-2.38,-0.94) circle (.05);
\fill[kDoc] (-2.69,-1.02) circle (.05);
\fill[kDoc] (-2.13,-0.77) circle (.05);
\fill[kDoc] (-2.44,-0.85) circle (.05);
\fill[kDoc] (-2.75,-0.91) circle (.05);
\fill[kDoc] (-2.18,-0.69) circle (.05);
\fill[kDoc] (-2.49,-0.75) circle (.05);
\fill[kDoc] (-2.81,-0.80) circle (.05);
\fill[kDoc] (-2.22,-0.60) circle (.05);
\fill[kDoc] (-2.54,-0.65) circle (.05);
\fill[kDoc] (-2.85,-0.69) circle (.05);
\fill[kDoc] (-2.25,-0.52) circle (.05);
\fill[kDoc] (-2.57,-0.55) circle (.05);
\fill[kDoc] (-2.89,-0.58) circle (.05);
\fill[kDoc] (-2.28,-0.43) circle (.05);
\fill[kDoc] (-2.61,-0.45) circle (.05);
\fill[kDoc] (-2.93,-0.47) circle (.05);
\fill[kDoc] (-2.31,-0.34) circle (.05);
\fill[kDoc] (-2.63,-0.35) circle (.05);
\fill[kDoc] (-2.95,-0.36) circle (.05);
\fill[kDoc] (-2.33,-0.25) circle (.05);
\fill[kDoc] (-2.65,-0.25) circle (.05);
\fill[kDoc] (-2.97,-0.24) circle (.05);
\fill[kDoc] (-2.34,-0.16) circle (.05);
\fill[kDoc] (-2.66,-0.15) circle (.05);
\fill[kDoc] (-2.99,-0.13) circle (.05);
\fill[kDoc] (-2.35,-0.07) circle (.05);
\fill[kDoc] (-2.67,-0.04) circle (.05);
\fill[kDoc] (-2.99,-0.01) circle (.05);
\fill[kDoc] (-2.35,0.02) circle (.05);
\fill[kDoc] (-2.67,0.06) circle (.05);
\fill[kDoc] (-2.99,0.11) circle (.05);
\fill[kDoc] (-2.35,0.11) circle (.05);
\fill[kDoc] (-2.66,0.16) circle (.05);
\fill[kDoc] (-2.98,0.22) circle (.05);
\fill[kDoc] (-2.34,0.20) circle (.05);
\fill[kDoc] (-2.65,0.27) circle (.05);
\fill[kDoc] (-2.96,0.34) circle (.05);
\fill[kDoc] (-2.32,0.29) circle (.05);
\fill[kDoc] (-2.63,0.37) circle (.05);
\fill[kDoc] (-2.93,0.45) circle (.05);
\fill[kDoc] (-2.30,0.38) circle (.05);
\fill[kDoc] (-2.60,0.47) circle (.05);
\fill[kDoc] (-2.90,0.56) circle (.05);
\fill[kDoc] (-2.27,0.47) circle (.05);
\fill[kDoc] (-2.57,0.57) circle (.05);
\fill[kDoc] (-2.86,0.68) circle (.05);
\fill[kDoc] (-2.24,0.56) circle (.05);
\fill[kDoc] (-2.53,0.67) circle (.05);
\fill[kDoc] (-2.81,0.79) circle (.05);
\fill[kDoc] (-2.20,0.65) circle (.05);
\fill[kDoc] (-2.48,0.77) circle (.05);
\fill[kDoc] (-2.76,0.89) circle (.05);
\fill[kDoc] (-2.15,0.73) circle (.05);
\fill[kDoc] (-2.43,0.86) circle (.05);
\fill[kDoc] (-2.70,1.00) circle (.05);
\fill[kDoc] (-2.11,0.81) circle (.05);
\fill[kDoc] (-2.37,0.96) circle (.05);
\fill[kDoc] (-2.63,1.10) circle (.05);
\fill[kDoc] (-2.05,0.89) circle (.05);
\fill[kDoc] (-2.31,1.05) circle (.05);
\fill[kDoc] (-2.56,1.20) circle (.05);
\fill[kDoc] (-1.99,0.97) circle (.05);
\fill[kConcept] (-2.24,1.13) circle (.05);
\fill[kConcept] (-2.48,1.30) circle (.05);
\fill[kConcept] (-1.93,1.05) circle (.05);
\fill[kConcept] (-2.16,1.22) circle (.05);
\fill[kConcept] (-2.39,1.40) circle (.05);
\fill[kConcept] (-1.86,1.12) circle (.05);
\fill[kConcept] (-2.08,1.30) circle (.05);
\fill[kConcept] (-2.30,1.49) circle (.05);
\fill[kConcept] (-1.78,1.19) circle (.05);
\fill[kConcept] (-2.00,1.38) circle (.05);
\fill[kConcept] (-2.20,1.58) circle (.05);
\fill[kConcept] (-1.71,1.26) circle (.05);
\fill[kConcept] (-1.91,1.46) circle (.05);
\fill[kConcept] (-2.10,1.66) circle (.05);
\fill[kConcept] (-1.62,1.33) circle (.05);
\fill[kConcept] (-1.81,1.53) circle (.05);
\fill[kConcept] (-1.99,1.74) circle (.05);
\fill[kConcept] (-1.54,1.39) circle (.05);
\fill[kConcept] (-1.71,1.60) circle (.05);
\fill[kConcept] (-1.88,1.81) circle (.05);
\fill[kConcept] (-1.45,1.44) circle (.05);
\fill[kWikidata] (-1.61,1.66) circle (.05);
\fill[kWikidata] (-1.76,1.88) circle (.05);
\fill[kWikidata] (-1.35,1.50) circle (.05);
\fill[kWikidata] (-1.50,1.72) circle (.05);
\fill[kWikidata] (-1.64,1.95) circle (.05);
\fill[kWikidata] (-1.26,1.55) circle (.05);
\fill[kWikidata] (-1.39,1.78) circle (.05);
\fill[kWikidata] (-1.51,2.01) circle (.05);
\fill[kWikidata] (-1.16,1.60) circle (.05);
\fill[kWikidata] (-1.27,1.83) circle (.05);
\fill[kWikidata] (-1.38,2.07) circle (.05);
\fill[kWikidata] (-1.05,1.64) circle (.05);
\fill[kWikidata] (-1.16,1.88) circle (.05);
\fill[kWikidata] (-1.25,2.12) circle (.05);
\fill[kWikidata] (-0.95,1.68) circle (.05);
\fill[kWikidata] (-1.03,1.92) circle (.05);
\fill[kWikidata] (-1.11,2.16) circle (.05);
\fill[kWikidata] (-0.84,1.71) circle (.05);
\fill[kTemplate] (-0.91,1.96) circle (.05);
\fill[kTemplate] (-0.97,2.21) circle (.05);
\fill[kTemplate] (-0.73,1.74) circle (.05);
\fill[kTemplate] (-0.78,1.99) circle (.05);
\fill[kTemplate] (-0.83,2.24) circle (.05);
\fill[kTemplate] (-0.62,1.77) circle (.05);
\fill[kTemplate] (-0.66,2.02) circle (.05);
\fill[kTemplate] (-0.69,2.27) circle (.05);
\fill[kStub] (-0.50,1.79) circle (.05);
\fill[kStub] (-0.53,2.04) circle (.05);
\fill[kStub] (-0.54,2.29) circle (.05);
\fill[kStub] (-0.39,1.81) circle (.05);
\fill[kStub] (-0.40,2.06) circle (.05);
\fill[kStub] (-0.40,2.31) circle (.05);
\fill[kConcept] (-0.27,1.82) circle (.05);
\fill[kConcept] (-0.27,2.07) circle (.05);
\fill[kConcept] (-0.25,2.32) circle (.05);
\fill[kConcept] (-0.16,1.83) circle (.05);
\fill[kConcept] (-0.13,2.08) circle (.05);
\fill[kConcept] (-0.10,2.33) circle (.05);
\fill[kConcept] (-0.04,1.83) circle (.05);
\fill[kFrame] (0,0) circle (.38); \node[font=\scriptsize\bfseries, text=white] at (0,0) {Doge};
\fill[kFrame] (4.0,1.90) circle (.07);\node[anchor=west, font=\scriptsize, text=gmInk] at (4.12,1.90) {156~frames};
\fill[kEvent] (4.0,1.56) circle (.07);\node[anchor=west, font=\scriptsize, text=gmInk] at (4.12,1.56) {64~events};
\fill[kDoc] (4.0,1.22) circle (.07);\node[anchor=west, font=\scriptsize, text=gmInk] at (4.12,1.22) {62~external pages};
\fill[kDoc] (4.0,0.88) circle (.07);\node[anchor=west, font=\scriptsize, text=gmInk] at (4.12,0.88) {37~images};
\fill[kConcept] (4.0,0.54) circle (.07);\node[anchor=west, font=\scriptsize, text=gmInk] at (4.12,0.54) {21~tags};
\fill[kWikidata] (4.0,0.20) circle (.07);\node[anchor=west, font=\scriptsize, text=gmInk] at (4.12,0.20) {18~Wikidata items};
\fill[kTemplate] (4.0,-0.14) circle (.07);\node[anchor=west, font=\scriptsize, text=gmInk] at (4.12,-0.14) {8~templates};
\fill[kStub] (4.0,-0.48) circle (.07);\node[anchor=west, font=\scriptsize, text=gmInk] at (4.12,-0.48) {6~frame stubs};
\fill[kConcept] (4.0,-0.82) circle (.07);\node[anchor=west, font=\scriptsize, text=gmInk] at (4.12,-0.82) {5~entry types};
\fill[kConcept] (4.0,-1.16) circle (.07);\node[anchor=west, font=\scriptsize, text=gmInk] at (4.12,-1.16) {1~region};
\fill[kConcept] (4.0,-1.50) circle (.07);\node[anchor=west, font=\scriptsize, text=gmInk] at (4.12,-1.50) {1~origin};
\node[anchor=west, font=\scriptsize\bfseries, text=gmInk] at (3.9,2.35) {379 edges};
\end{tikzpicture}
}
  \note{
  \begin{itemize}
    \item The previous picture had names; this one has everything: 379 edges touch Doge.
    \item 156 other entries (links between pages, siblings, events pointing at other memes), 64
      events and event links, 62 external pages (its references), 37 images, 21 tags, 18
      Wikidata items, 8 templates \dots
    \item One meme. The whole graph has 24,291 of them.
  \end{itemize}}
\end{frame}

\begin{frame}[fragile]{The same facts, in two languages}
  \begin{columns}[t]
    \column{.54\textwidth}
      {\scriptsize\bfseries\color{gm@col@graph}RDF --- for the Semantic Web, queried with SPARQL}\par\vskip2mm
      {\fontsize{6.6}{8.4}\selectfont\ttfamily
      kym:doge \textcolor{gmInk2}{a} m4s:MediaFrame, kym:Meme ;\\
      \hspace*{4mm}m4s:year 2010 ;\\
      \hspace*{4mm}m4s:fromTags wd:Q144 ;\\
      \hspace*{4mm}skos:broader kym:interior-monologue-captioning ;\\
      \hspace*{4mm}rdfs:seeAlso kym:animals-talking-in-all-caps ;\\
      \hspace*{4mm}mk:hasTemplate mk:template/8072285 .\\[2mm]
      \textcolor{gmInk2}{<<} kym:doge m4s:fromTags wd:Q144 \textcolor{gmInk2}{>>}\\
      \hspace*{4mm}mk:mentionText "dog" ; mk:linkScore 0.616 .}
    \column{.42\textwidth}
      {\scriptsize\bfseries\color{gm@col@graph}a property graph --- Neo4j, queried with Cypher}\par\vskip2mm
      {\fontsize{6.6}{8.4}\selectfont\ttfamily
      MATCH (f:Frame \{label: 'Doge'\})\\
      \hspace*{4mm}-[r:fromTags]->\\
      \hspace*{4mm}(e:WikidataEntity \{qid: 'Q144'\})\\
      RETURN f.year, r.mention\_texts,\\
      \hspace*{4mm}r.link\_scores}\par\vskip5mm
      \begin{tikzpicture}[baseline=(current bounding box.north)]
        \node[fill=gmTint, rounded corners=2mm, inner sep=3mm, text width=4.6cm, align=left] {
          \gmwhy{IMKG's words, kept}\par\vskip1mm\scriptsize
          m4s:, kym:, skos:broader, rdfs:seeAlso --- word for word. MemeAtlas's own terms (mk:) only
          where IMKG has none.};
      \end{tikzpicture}
  \end{columns}
  \note{
  \begin{itemize}
    \item The graph is published twice, from the same data: as RDF (W3C standard, queried with
      SPARQL, in Fuseki) and as a property graph (Neo4j, Cypher).
    \item Left: Doge in RDF. m4s: and kym: are IMKG's own vocabulary; series are skos:broader;
      siblings are rdfs:seeAlso --- exactly what IMKG emits, so IMKG's queries still run. mk: is
      MemeAtlas's namespace, for what IMKG has no word for (templates, events, curation\dots).
    \item The double angle brackets are RDF-star: a statement about a statement. ``Doge's tags
      mention dog'' carries the text and the score of that link.
    \item Right: the same fact in Cypher. Property-graph edges carry those details as properties.
  \end{itemize}}
\end{frame}

\begin{frame}{KG \NVersion: \NNodes\ nodes, \NEdges\ edges}
  \begin{columns}[c]
    \column{.5\textwidth}
      % generated by scripts/figures.py — do not edit
\begin{tikzpicture}[x=1cm, y=-.42cm]
\node[anchor=east, font=\scriptsize] at (-.15,0) {Wikidata items};
\fill[kWikidata] (0,-0.3) rectangle (6.200,0.3);
\node[anchor=west, font=\tiny, text=gmInk2] at (6.280,0) {272,304};
\node[anchor=east, font=\scriptsize] at (-.15,1) {external pages};
\fill[kDoc] (0,0.7) rectangle (5.852,1.3);
\node[anchor=west, font=\tiny, text=gmInk2] at (5.932,1) {257,032};
\node[anchor=east, font=\scriptsize] at (-.15,2) {images};
\fill[kDoc] (0,1.7) rectangle (4.607,2.3);
\node[anchor=west, font=\tiny, text=gmInk2] at (4.687,2) {202,347};
\node[anchor=east, font=\scriptsize] at (-.15,3) {events};
\fill[kEvent] (0,2.7) rectangle (3.131,3.3);
\node[anchor=west, font=\tiny, text=gmInk2] at (3.211,3) {137,496};
\node[anchor=east, font=\scriptsize] at (-.15,4) {tags};
\fill[kConcept] (0,3.7) rectangle (2.320,4.3);
\node[anchor=west, font=\tiny, text=gmInk2] at (2.400,4) {101,882};
\node[anchor=east, font=\scriptsize] at (-.15,5) {templates};
\fill[kTemplate] (0,4.7) rectangle (0.612,5.3);
\node[anchor=west, font=\tiny, text=gmInk2] at (0.692,5) {26,868};
\node[anchor=east, font=\scriptsize] at (-.15,6) {frames};
\fill[kFrame] (0,5.7) rectangle (0.535,6.3);
\node[anchor=west, font=\tiny, text=gmInk2] at (0.615,6) {23,477};
\node[anchor=east, font=\scriptsize] at (-.15,7) {frame stubs};
\fill[kStub] (0,6.7) rectangle (0.174,7.3);
\node[anchor=west, font=\tiny, text=gmInk2] at (0.254,7) {7,629};
\node[anchor=east, font=\scriptsize] at (-.15,8) {origins};
\fill[kConcept] (0,7.7) rectangle (0.130,8.3);
\node[anchor=west, font=\tiny, text=gmInk2] at (0.210,8) {5,703};
\node[anchor=east, font=\scriptsize] at (-.15,9) {entry types};
\fill[kConcept] (0,8.7) rectangle (0.030,9.3);
\node[anchor=west, font=\tiny, text=gmInk2] at (0.110,9) {119};
\node[anchor=east, font=\scriptsize] at (-.15,10) {regions};
\fill[kConcept] (0,9.7) rectangle (0.030,10.3);
\node[anchor=west, font=\tiny, text=gmInk2] at (0.110,10) {110};
\node[anchor=east, font=\scriptsize] at (-.15,11) {badges};
\fill[kConcept] (0,10.7) rectangle (0.030,11.3);
\node[anchor=west, font=\tiny, text=gmInk2] at (0.110,11) {1};
\end{tikzpicture}

    \column{.46\textwidth}
      \gmbig{\NTriplesShort}{RDF triples}\par\vskip4mm
      \gmbig{\NRelTypes}{kinds of relation}\par\vskip4mm
      {\small\color{gmInk2}published \NPublished\ as build {\ttfamily\scriptsize\NBuildId}}
  \end{columns}
  \note{
  \begin{itemize}
    \item The live graph, version 7.0.0, published 2 October: 795,711 nodes, 2,564,089 edges,
      8.79 million RDF triples, 26 relation types.
    \item Bars are node kinds, in the same colours as the pictures: external pages and images are
      the most numerous (every reference and image of every page), then events, tags, Wikidata
      items, templates, and the 24,291 frames themselves.
  \end{itemize}}
\end{frame}

\begin{frame}{The graph of the graph}
  \centering
  \resizebox{!}{6.3cm}{% generated by scripts/figures.py — do not edit
\begin{tikzpicture}[x=1cm, y=1cm, lab/.style={font=\fontsize{8}{9}\selectfont, text=gmInk2, fill=gmSurface, inner sep=.7pt}]
\node[circle, fill=kFrame, minimum size=0.97cm, inner sep=0pt] (k-frame) at (4.3,3.05) {};
\node[font=\fontsize{8.5}{9.5}\selectfont, text=gmInk, align=center, below=.5mm of k-frame, fill=gmSurface, inner sep=.6pt] {\textbf{frame} 24,291\\{\color{gmInk2}$\circlearrowleft$ siblings, links 856k}};
\node[circle, fill=kEvent, minimum size=1.77cm, inner sep=0pt] (k-event) at (8.5,3.05) {};
\node[font=\fontsize{8.5}{9.5}\selectfont, text=gmInk, align=center, below=.5mm of k-event, fill=gmSurface, inner sep=.6pt] {\textbf{event} 142,787\\{\color{gmInk2}$\circlearrowleft$ nextInStory 134k}};
\node[circle, fill=kDoc, minimum size=2.24cm, inner sep=0pt] (k-external_ref) at (11.7,5.5) {};
\node[font=\fontsize{8.5}{9.5}\selectfont, text=gmInk, align=center, below=.5mm of k-external_ref, fill=gmSurface, inner sep=.6pt] {\textbf{external page} 257,031};
\node[circle, fill=kDoc, minimum size=2.03cm, inner sep=0pt] (k-image) at (11.7,0.6) {};
\node[font=\fontsize{8.5}{9.5}\selectfont, text=gmInk, align=center, below=.5mm of k-image, fill=gmSurface, inner sep=.6pt] {\textbf{image} 202,390};
\node[circle, fill=kWikidata, minimum size=1.00cm, inner sep=0pt] (k-wikidata_entity) at (0.9,5.75) {};
\node[font=\fontsize{8.5}{9.5}\selectfont, text=gmInk, align=center, below=.5mm of k-wikidata_entity, fill=gmSurface, inner sep=.6pt] {\textbf{Wikidata item} 26,897};
\node[circle, fill=kTemplate, minimum size=0.99cm, inner sep=0pt] (k-template) at (6.6,0.0) {};
\node[font=\fontsize{8.5}{9.5}\selectfont, text=gmInk, align=center, below=.5mm of k-template, fill=gmSurface, inner sep=.6pt] {\textbf{template} 26,868};
\node[circle, fill=kStub, minimum size=0.72cm, inner sep=0pt] (k-frame_stub) at (6.0,6.05) {};
\node[font=\fontsize{8.5}{9.5}\selectfont, text=gmInk, align=center, below=.5mm of k-frame_stub, fill=gmSurface, inner sep=.6pt] {\textbf{frame stub} 7,631};
\node[circle, fill=kConcept, minimum size=1.56cm, inner sep=0pt] (k-tag_concept) at (0.7,2.75) {};
\node[font=\fontsize{8.5}{9.5}\selectfont, text=gmInk, align=center, below=.5mm of k-tag_concept, fill=gmSurface, inner sep=.6pt] {\textbf{tag} 101,883\\{\color{gmInk2}$\circlearrowleft$ coOccursWith 5k}};
\node[circle, fill=kConcept, minimum size=0.44cm, inner sep=0pt] (k-entry_type_concept) at (1.3,0.35) {};
\node[font=\fontsize{8.5}{9.5}\selectfont, text=gmInk, align=center, below=.5mm of k-entry_type_concept, fill=gmSurface, inner sep=.6pt] {\textbf{entry type} 119};
\node[circle, fill=kConcept, minimum size=0.67cm, inner sep=0pt] (k-origin_concept) at (3.4,-0.45) {};
\node[font=\fontsize{8.5}{9.5}\selectfont, text=gmInk, align=center, below=.5mm of k-origin_concept, fill=gmSurface, inner sep=.6pt] {\textbf{origin} 5,703};
\node[circle, fill=kConcept, minimum size=0.44cm, inner sep=0pt] (k-region_concept) at (3.2,6.05) {};
\node[font=\fontsize{8.5}{9.5}\selectfont, text=gmInk, align=center, below=.5mm of k-region_concept, fill=gmSurface, inner sep=.6pt] {\textbf{region} 110};
\begin{pgfonlayer}{background}
\draw[gmBase, line width=2.56pt, -{Stealth[length=2.9mm]}] (k-frame) -- node[lab, sloped, above, pos=0.38] {citesExternal 244k} (k-external_ref);
\draw[gmBase, line width=2.54pt, -{Stealth[length=2.9mm]}] (k-frame) -- node[lab, sloped, above, pos=0.5] {hasTag 226k} (k-tag_concept);
\draw[gmBase, line width=2.50pt, -{Stealth[length=2.9mm]}] (k-frame) to[bend right=10] node[lab, sloped, above, pos=0.3] {hasImage 193k} (k-image);
\draw[gmBase, line width=2.49pt, -{Stealth[length=2.8mm]}] (k-frame) -- node[lab, sloped, above, pos=0.4] {about, tags, title 181k} (k-wikidata_entity);
\draw[gmBase, line width=2.42pt, -{Stealth[length=2.8mm]}] (k-frame) to[bend left=14] node[lab, sloped, above, pos=0.5] {hasEvent 143k} (k-event);
\draw[gmBase, line width=2.41pt, -{Stealth[length=2.8mm]}] (k-event) -- node[lab, sloped, above, pos=0.38] {citations, embeds 138k} (k-external_ref);
\draw[gmBase, line width=2.34pt, -{Stealth[length=2.8mm]}] (k-event) -- node[lab, sloped, above, pos=0.45] {eventImage 104k} (k-image);
\draw[gmBase, line width=2.23pt, -{Stealth[length=2.7mm]}] (k-event) to[bend left=14] node[lab, sloped, above, pos=0.5] {eventLink 67k} (k-frame);
\draw[gmBase, line width=2.19pt, -{Stealth[length=2.7mm]}] (k-frame) -- node[lab, sloped, above, pos=0.45] {citesMediaFrame 58k} (k-frame_stub);
\draw[gmBase, line width=2.12pt, -{Stealth[length=2.7mm]}] (k-template) to[bend left=26] node[lab, sloped, above, pos=0.7] {fromImage 45k} (k-wikidata_entity);
\draw[gmBase, line width=2.05pt, -{Stealth[length=2.6mm]}] (k-frame) -- node[lab, sloped, above, pos=0.55] {hasEntryType 34k} (k-entry_type_concept);
\draw[gmBase, line width=2.02pt, -{Stealth[length=2.6mm]}] (k-frame) -- node[lab, sloped, above, pos=0.5] {hasTemplate 30k} (k-template);
\draw[gmBase, line width=1.99pt, -{Stealth[length=2.6mm]}] (k-template) to[bend right=24] node[lab, sloped, above, pos=0.55] {imgflipPage 27k} (k-external_ref);
\draw[gmBase, line width=1.97pt, -{Stealth[length=2.6mm]}] (k-template) -- node[lab, sloped, above, pos=0.5] {templateImage 25k} (k-image);
\draw[gmBase, line width=1.96pt, -{Stealth[length=2.6mm]}] (k-frame) -- (k-origin_concept);
\draw[gmBase, line width=1.94pt, -{Stealth[length=2.6mm]}] (k-event) to[bend right=12] (k-frame_stub);
\draw[gmBase, line width=1.58pt, -{Stealth[length=2.4mm]}] (k-frame) -- (k-region_concept);
\end{pgfonlayer}
\end{tikzpicture}
}
  \note{
  \begin{itemize}
    \item Every kind of node, every kind of edge between kinds, sized by how many there are
      (area: nodes; width: edges, on a log scale).
    \item The frame is the hub. Biggest edges: siblings (679k, derived from the series), links to
      external pages (244k), tags (226k), images (193k), events (143k) and the story chain between
      events (124k).
    \item This is the schema of the graph, measured, not designed on a whiteboard.
  \end{itemize}}
\end{frame}

\begin{frame}{One build, four stores, nothing published unchecked}
  \centering
  \begin{tikzpicture}[x=1cm, y=1cm,
      st/.style={rounded corners=2mm, draw=gm@col@graph, line width=1pt, fill=gmSurface, minimum width=2.5cm,
                 minimum height=1.25cm, align=center, font=\small},
      go/.style={-{Stealth[length=2.2mm]}, line width=.9pt, draw=gmBase}]
    \node[st, fill=gm@col@graph, text=white, minimum width=2.6cm] (b) at (0,0) {\bfseries one build\\{\scriptsize build.py}};
    \node[st] (m) at (4.2,1.55) {Mongo\\{\scriptsize the authority}};
    \node[st] (f) at (4.2,0) {files\\{\scriptsize graph.nt, CSV}};
    \node[st] (u) at (4.2,-1.55) {Fuseki\\{\scriptsize SPARQL}};
    \node[st] (n) at (7.3,-1.55) {Neo4j\\{\scriptsize Cypher}};
    \draw[go] (b) -- (m); \draw[go] (b) -- (f); \draw[go] (f) -- (u); \draw[go] (b.south) |- (n.west);
    \node[st, draw=kTemplate, minimum width=3.3cm] (v) at (10.9,1.1) {verify\\{\scriptsize every count agrees}};
    \node[st, draw=kTemplate, minimum width=3.3cm] (p) at (10.9,-1.0) {publish\\{\scriptsize flip every pointer}};
    \draw[go] (m.east) -- ++(1.2,0) |- (v.west);
    \draw[go] (n.east) -- ++(.5,0) |- (v.west);
    \draw[go] (v) -- (p);
    \node[font=\scriptsize, text=gmInk2, align=left, anchor=north west, text width=13.4cm] at (-1.3,-2.55)
      {and once a week the RDF is re-derived another way, through the mapping, and compared:
       \NRmlTriples\ against \NTriples\ --- equal but for \NTriplesDiff\ provenance triples.};
  \end{tikzpicture}
  \note{
  \begin{itemize}
    \item kym\_kg: build.py is the only producer of nodes and edges. One stream feeds every
      representation: Mongo (the authority), files (N-Triples, CSVs for RML, view CSVs), Fuseki
      (RDF) and Neo4j (property graph).
    \item Builds are generational: each has an id, stores keep the live build and the one before,
      for rollback. Publishing flips pointers --- followers first, the authority last --- and reads
      them back. A torn publish fails loudly.
    \item Before publishing, verify checks that Mongo, the files, Fuseki and Neo4j hold the same
      counts. If not, nothing moves.
    \item Independently, kym\_kg\_validate re-derives the RDF from the CSVs through the YARRRML
      mapping (morph-kgc) and diffs it with graph.nt: 8,790,365 vs 8,790,369 --- the 4 extra are
      the build's own provenance triples.
  \end{itemize}}
\end{frame}

\begin{frame}{A graph with a version number}
  \begin{columns}[t]
    \column{.5\textwidth}
      \begin{itemize}\setlength{\itemsep}{2.6mm}
        \item Every input carries a \textbf{stamp}: parser, prompt, model, lexicon, taxonomy versions
        \item The build compares stamps: \textbf{nothing changed, nothing rebuilt}
        \item The graph has a \textbf{semantic version}: 7.0.0 renamed a relation, so it is a
          major release; 7.1.0 only adds
        \item \FDifferences\ deliberate differences from IMKG, each written down
      \end{itemize}
    \column{.44\textwidth}
      {\small\color{gmInk2}from the first graph to today}\par\vskip2mm
      {\small
      \begin{tabular}{@{}rl@{}}
        \textbf{4.0} & a node only for what things share\\
        \textbf{5.0} & origins, badges, folded tags\\
        \textbf{6.0} & events\\
        \textbf{6.1} & Wikidata links\\
        \textbf{6.4} & templates and their images\\
        \textbf{6.5} & curated links\\
        \textbf{7.0} & series siblings; citesMediaFrame\\
        \textbf{7.1} & frame images; Wikidata facts\\
      \end{tabular}}
  \end{columns}
  \note{
  \begin{itemize}
    \item Staleness: the build records the version of everything that could change the output
      (parser, taxonomies, extraction prompt and model, lexicon, curation, templates\dots). The
      monthly run compares them with the published build; if nothing moved, it stops there.
    \item Versions follow semantic versioning: a renamed term breaks queries, so 7.0.0 is major
      (relatesToMeme became citesMediaFrame); a new layer only adds, so 7.1.0 is minor.
    \item MODEL.md documents the crosswalk to IMKG and 13 deliberate differences --- e.g. siblings
      derived from the series instead of KYM's truncated ``related entries'' box, or the Wikidata
      entity IRI instead of the wiki page URL.
    \item A version worth telling: 4.0.0 decided what deserves to be a node --- only what things
      share. Sections, links and references had been nodes; half of them were detours. The graph
      went from 985,481 to 476,794 nodes and 1.74M to 856k edges, with no data lost.
  \end{itemize}}
\end{frame}

\gmzoomout{graph}
